\section{Proofs for \texorpdfstring{Section~\ref{sec:sca}}{scale}}\label{app:sca}
% src: proofs of L4, B3 Sections 1, 2, 5, B5 Section 1, B1-P3, B1-P6, L6-P5(c), as cited in each proof; prop:sca-records-length(b) proved by a finite-dimensional argument written in this editing round

We prove the results of \cref{sec:sca} in order. Block norms are spectral norms, $\Phi(t,s)=W(t-1)\cdots W(s)$, and we use repeatedly that block row sums of a product of block matrices are at most the products of the block row sums, and likewise for column sums: $\sum_j\Vert(AB)_{ij}\Vert\le\sum_k\Vert A_{ik}\Vert\max_k\sum_j\Vert B_{kj}\Vert$.

\begin{proof}[Proof of \cref{lem:sca-relations}]
(a) The row-sum inequality and induction give $\sum_j\Vert\Phi_{ij}(t,s)\Vert\le\lambda^{t-s}$. Write $e_i(t)=\sum_k\Phi_{ik}(t,0)e_k(0)+\sum_{u<t}\sum_k\Phi_{ik}(t,u+1)\zeta_k(u)$. Under \textup{(Lin)} the noise terms are centered, independent of one another and of the first sum. With the weights $\Vert\Phi_{ik}(t,0)\Vert$ the Cauchy--Schwarz inequality gives
\begin{equation*}
\E\Bigl\Vert\sum_k\Phi_{ik}(t,0)e_k(0)\Bigr\Vert^2\le\Bigl(\sum_k\Vert\Phi_{ik}(t,0)\Vert\Bigr)\sum_k\Vert\Phi_{ik}(t,0)\Vert\E\Vert e_k(0)\Vert^2\le\lambda^{2t}\sup_k\E\Vert e_k(0)\Vert^2,
\end{equation*}
and for block-diagonal noise, with $k_{\max}$ the largest block dimension,
\begin{equation*}
\E\Bigl\Vert\sum_k\Phi_{ik}(t,u+1)\zeta_k(u)\Bigr\Vert^2=\sum_k\tr\bigl(\Phi_{ik}Q^\zeta_k(u)\Phi_{ik}^\top\bigr)\le qk_{\max}\Bigl(\sum_k\Vert\Phi_{ik}(t,u+1)\Vert\Bigr)^2\le qk_{\max}\lambda^{2(t-1-u)}.
\end{equation*}
Summing, $\E\Vert e_i(t)\Vert^2\le\lambda^{2t}\sup_k\E\Vert e_k(0)\Vert^2+qk_{\max}/(1-\lambda^2)$, which is finite and uniform in $i$ and $t$ whenever $\sup_k\E\Vert e_k(0)\Vert^2<\infty$. For the uniform form of \textup{(M)}, write $\tilde e_i(t)=m_i(t)-\eta_t$ for the errors themselves, which obey $\tilde e(t+1)=W(t)\tilde e(t)+\zeta(t)+B^{\mathrm{mis}}\eta_t-\one\otimes w_t$ (\cref{lem:fw-common-mode}). Then $\tilde e_i(t)$ is the sum of the initial term $\sum_k\Phi_{ik}(t,0)\tilde e_k(0)$, the noise terms of the second display and the input terms $\sum_k\Phi_{ik}(t,u+1)(B^{\mathrm{mis}}_k\eta_u-w_u)$, $u<t$. Because the schedule is independent of the initial errors, the first display holds with $\tilde e_k(0)$ in place of $e_k(0)$, so the initial term is at most $\lambda^t\sup_k(\E\Vert m_k(0)-\eta_0\Vert^2)^{1/2}$ in root mean square. The noise terms are centered and uncorrelated across rounds, so together they contribute at most $(qk_{\max}/(1-\lambda^2))^{1/2}$. Pathwise $\Vert\sum_k\Phi_{ik}(t,u+1)(B^{\mathrm{mis}}_k\eta_u-w_u)\Vert\le\lambda^{t-1-u}(b\Vert\eta_u\Vert+\Vert w_u\Vert)$ with $b=\sup_k\Vert B^{\mathrm{mis}}_k\Vert$, so by Minkowski's inequality the input terms contribute at most $\bigl(b\sup_u(\E\Vert\eta_u\Vert^2)^{1/2}+\sup_u(\E\Vert w_u\Vert^2)^{1/2}\bigr)/(1-\lambda)$, where the first summand is absent when $B^{\mathrm{mis}}=0$. Minkowski's inequality applied to the three contributions bounds $(\E\Vert\tilde e_i(t)\Vert^2)^{1/2}$ uniformly in $i$ and $t$; this is the uniform form of \textup{(M)}. For time-invariant $W$ and identically distributed $\zeta(t)$ the series $\sum_{n\ge0}W^n\zeta(t-1-n)$ converges in mean square, blockwise with $\E\Vert(\cdot)_i\Vert^2\le qk_{\max}/(1-\lambda^2)$, and its law does not depend on $t$, which defines the stationary law; for Gaussian noise the law of $\zeta(t)$ is fixed by $Q^\zeta$. For a schedule that depends on the states, on $N$ agents, write the recursion of the errors as $e(t+1)=W(t)e(t)+u(t)$, where $u(t)$ collects all inputs, local noise, world innovation and any mistranslation term, with $\sup_{k,t}\E\Vert u_k(t)\Vert^2=U<\infty$. Every realized $W(t)$ has block row sums at most $\lambda$, so pathwise $\max_i\Vert e_i(t+1)\Vert\le\lambda\max_k\Vert e_k(t)\Vert+\max_k\Vert u_k(t)\Vert$, and by induction $\Vert e_i(t)\Vert\le\lambda^t\max_k\Vert e_k(0)\Vert+\sum_{s<t}\lambda^{t-1-s}\max_k\Vert u_k(s)\Vert$. Minkowski's inequality with $\E\max_k\Vert u_k(s)\Vert^2\le\sum_k\E\Vert u_k(s)\Vert^2\le NU$ gives $(\E\Vert e_i(t)\Vert^2)^{1/2}\le(\E\max_k\Vert e_k(0)\Vert^2)^{1/2}+(NU)^{1/2}/(1-\lambda)$, finite and uniform in $i$ and $t$. No independence between the schedule and the inputs is used.

(b) Let $b_{ij}=\Vert W_{ij}\Vert$. For a vector with blocks $x_j$, $\Vert(Wx)_i\Vert\le\sum_jb_{ij}\Vert x_j\Vert$, so $\Vert W\Vert_2\le\Vert(b_{ij})\Vert_2\le(\max_i\sum_jb_{ij}\cdot\max_j\sum_ib_{ij})^{1/2}$ by the Schur test. Hence $\Vert W^k\Vert_2\le(\lambda_r\lambda_c)^{k/2}$, which is \textup{(C$_\rho$)} with $K_A=1$ and $f=(\lambda_r\lambda_c)^{1/2}$.

(c) If $\rho(W)<1$, Gelfand's formula for the bounded operator $W$ on $\ell^2(V)$ gives $\Vert W^n\Vert_2\le K_Af^n$ for some $f<1$, and since $Q^\zeta$ is bounded the series $\sum_nW^nQ^\zeta W^{n\top}$ converges in operator norm to the covariance of the stationary solution $\sum_{n\ge0}W^n\zeta(t-1-n)$, whose blocks are therefore finite. Conversely, let the population be finite. A stationary covariance satisfies $\Sigma=W\Sigma W^\top+Q^\zeta$. For a left eigenvector $y^\ast W=\Lambda y^\ast$ this gives $(1-|\Lambda|^2)y^\ast\Sigma y=y^\ast Q^\zeta y>0$, so $|\Lambda|<1$ for every eigenvalue $\Lambda$. The argument needs an eigenvector, which a bounded operator on an infinite population need not have at its spectral radius.

For scalar agents on $\Z^D$ under \textup{(TI)}, $W^m$ is the convolution with symbol $\hat W^m$, so by Parseval's identity $\Var\bigl((W^m\zeta(u))_0\bigr)=q\sum_x|(W^m)_{0x}|^2=q(2\pi)^{-D}\int|\hat W(p)|^{2m}\mathrm dp$. If \eqref{eq:sca-fourier-moment} holds, then $|\hat W|<1$ almost everywhere and monotone convergence gives $\sum_{m\ge0}(2\pi)^{-D}\int|\hat W|^{2m}\mathrm dp=(2\pi)^{-D}\int(1-|\hat W|^2)^{-1}\mathrm dp<\infty$, so $\sum_{n\ge0}W^n\zeta(t-1-n)$ converges in mean square in every coordinate and defines a stationary law with the stated variance. Conversely, under a stationary law with finite second moments, $e(t)=W^ne(t-n)+\sum_{m<n}W^m\zeta(t-1-m)$ with $e(t-n)$ independent of the later noises, so $\Var(e_0(t))\ge q\sum_{m<n}(2\pi)^{-D}\int|\hat W|^{2m}\mathrm dp$ for every $n$; letting $n\to\infty$, monotone convergence gives \eqref{eq:sca-fourier-moment}. For the lazy walk on $\Z^3$, $\hat W(p)=\tfrac14+\tfrac14\sum_d\cos p_d\in[-\tfrac12,1]$ attains one at $p=0$, so $\rho(W)=1$, while $1-\hat W(p)^2\ge\tfrac12(1-\hat W(p))$ and $1-\hat W(p)\ge|p|^2/(2\pi^2)$ on $[-\pi,\pi]^3$, which makes the integrand at most $4\pi^2|p|^{-2}$, integrable in three dimensions. With $G(E)=\int_0^\infty e^{-Es}I_0(s)^3\mathrm ds$, the lattice Green function of the simple cubic lattice, partial fractions give the mean of the integrand as $2(G(3)+G(5))=1.4395$. On $\Z$ and $\Z^2$ a nearest-neighbor walk kernel with unit row sums has $1-\hat W(p)^2\le2(1-\hat W(p))\le C|p|^2$ near $p=0$, and $|p|^{-2}$ is not integrable there.
\end{proof}

\begin{proof}[Proof of \cref{prop:sca-sufficient}]
(a) \Cref{rem:fw-convex-gains} gives $W(t)\ge0$ with row sums $F(1-k^s_i(t))$, hence the bound on $\lambda$, and \cref{lem:sca-relations}(a) gives \textup{(M)}, through its pathwise bound when the schedule depends on the states. A sender attended by $L$ receivers has column sum at least $F\sum_ik_{ij}$ over them, which grows with $L$; \cref{cor:sca-size}(e) shows the consequence. (b) \Cref{prop:min-common-code}(b) gives $\rho(W)\le f$ and $\Vert W^k\Vert_2\le K_{\tilde M}f^k$, which is \textup{(C$_\rho$)}, and \cref{lem:sca-relations}(c) gives \textup{(M)}; the instances of \cref{chk:min-common-code} with block row sum above one show that \textup{(C$_r$)} can fail. (c) is the final assertion of \cref{prop:min-grounding-moment}, which applies when $\varpi_g<1$, and (d) is \cref{prop:sca-routing}(d).
\end{proof}

\begin{proof}[Proof of \cref{thm:sca-lightcone}]
(a) By induction on $t-s$, the case $t=s$ being trivial. Suppose the claim holds at $t$. The state $\mu_i(t+1)=\mathfrak U^{(t)}_i(\mu_{\mathcal G[i]}(t),\zeta_i(t))$ depends on $\mu_k(t)$ for $k\in\mathcal G[i]$, each a function of states $\mu_l(s)$ with $d(k,l)\le t-s$ and of noises $\zeta_l(u)$ with $d(k,l)\le t-1-u$. Since $d(i,k)\le1$, these have $d(i,l)\le t+1-s$ and $d(i,l)\le(t+1)-1-u$, and the new noise $\zeta_i(t)$ has $d(i,i)=0=(t+1)-1-t$. If $d(i,j)>t-s$, $\mu_i(t)$ does not depend on $\mu_j(s)$, so $J_{ij}(t,s)=0$.

(b) By (a) with $s=0$, $\mu_i(t)$ is a function of the family $\mathfrak F_i=\{\mu_k(0):d(i,k)\le t\}\cup\{\zeta_k(u):u<t,\ d(i,k)\le t-1-u\}$, and $\mu_j(t)$ of $\mathfrak F_j$. A common initial state would give $d(i,j)\le2t$, and a common $\zeta_k(u)$ would give $d(i,j)\le2(t-1-u)<2t$. So if $d(i,j)>2t$ the families are disjoint, and given the world path their members are independent; functions of disjoint independent families are independent. For $\mu_i(t+h)$ the family has radius $t+h$ for initial states and $t+h-1-u$ for $\zeta_k(u)$; a common initial state needs $d(i,j)\le2t+h$, and a common $\zeta_k(u)$ needs $u<t$ and then $d(i,j)\le2t+h-2-2u$. So $d(i,j)>2t+h$ suffices.

(c) The chain rule gives $J_{ij}(t+1,s)=\sum_{k\in\mathcal G[i]}\partial_k\mathfrak U^{(t)}_iJ_{kj}(t,s)$, so $\Vert J_{ij}(t+1,s)\Vert\le\sum_k\bar J_{ik}\Vert J_{kj}(t,s)\Vert$, and induction from $J(s,s)=I$ gives $(\bar J^{t-s})_{ij}$. Summing over $j$ uses $\sum_j(\bar J^n)_{ij}\le g_{\mathrm{tot}}^n$. If $(\bar J^n)_{kj}\le e^{-\varphi d(k,j)}g_\varphi^n$, then, since $-d(k,j)\le d(i,k)-d(i,j)$, $(\bar J^{n+1})_{ij}\le\sum_k\bar J_{ik}e^{-\varphi d(k,j)}g_\varphi^n\le e^{-\varphi d(i,j)}g_\varphi^{n+1}$.
\end{proof}

\begin{proof}[Proof of \cref{prop:sca-poisson}]
Fix the clocks and the noises; almost surely $[0,t]$ contains finitely many events. At an event of $l$ the Jacobian of the event map has row $l$ equal to $(\partial_k\mathfrak U_l)_k$ and identity rows elsewhere. Let $x(0)$ be the $j$-th coordinate vector and, at each event of $l$, set $x_l\leftarrow\bar J_{ll}x_l+\sum_{k\neq l}\bar J_{lk}x_k$, the other coordinates unchanged; by induction over events, $\Vert J_{kj}(s)\Vert\le x_k(s)$ for all $k$ and $s$. Expanding the recursion, $x_i(t)$ is a sum over walks $(i_0=j,i_1,\dots,i_n=i)$ with $i_m\neq i_{m-1}$ and over increasing times $t_1<\cdots<t_n$, $t_m$ an event of $i_m$, of $\prod_{m=1}^n\bar J_{i_mi_{m-1}}\prod_{m=0}^n\bar J_{i_mi_m}^{N_m}$, where $N_m$ counts the other events of $i_m$ in $(t_m,t_{m+1})$, $t_0=0$, $t_{n+1}=t$. By the multivariate Mecke formula for independent Poisson processes, the expected sum over increasing tuples is $\nu^n\int_{0<t_1<\cdots<t_n<t}\E[\prod_m\bar J_{i_mi_m}^{N_m}]\mathrm{d}t_1\cdots\mathrm{d}t_n$, with the $N_m$ now counts of the original processes over disjoint intervals, independent and Poisson with means $\nu(t_{m+1}-t_m)$. Hence $\E\prod_m\bar J_{i_mi_m}^{N_m}=\prod_me^{\nu(t_{m+1}-t_m)(\bar J_{i_mi_m}-1)}\le e^{-(1-\bar J_{\mathrm{diag}})\nu t}$. The time integral is $t^n/n!$, walks from $j$ to $i$ have $n\ge d$, and the sum of $\prod_m\bar J_{i_mi_{m-1}}$ over walks of length $n$ is $(\bar J_{\mathrm{off}}^n)_{ij}\le g^n$. This proves \eqref{eq:sca-poisson}. The explicit forms follow from $\sum_{n\ge d}x^n/n!\le(x^d/d!)\sum_{m\ge0}(x/(d+1))^m$ with $d!\ge\sqrt{2\pi d}(d/e)^d$, and from $\sum_{n\ge d}x^n/n!\le e^{-\varphi d}e^{xe^\varphi}$ with $x=g\nu t$. For $v>v_{\mathrm{LR}}$ there is $\varphi>0$ with $c_\varphi=\nu(ge^\varphi+\bar J_{\mathrm{diag}}-1)/\varphi<v$, and along $d\ge vt$ the second form is at most $e^{-\varphi(v-c_\varphi)t}$; when $g+\bar J_{\mathrm{diag}}<1$ the ratio $c_\varphi$ tends to $-\infty$ as $\varphi\downarrow0$, so every $v>0$ qualifies. Without smoothness, $\mu_i(t)$ can depend on $\mu_j(0)$ only through a walk from $j$ to $i$ whose steps occur at increasing event times; the expected number of such chains is $\sum_{n\ge d}N_n(i,j)(\nu t)^n/n!$, with $N_n(i,j)\le\Delta^n$ walks of length $n$. For linear scalar updates with nonnegative coefficients $A$ and $\diag(A)=\bar J_{\mathrm{diag}}I$, an event of agent $l$ in $[s,s+\mathrm{d}s)$ replaces row $l$ of $J$ by $A_{l\cdot}J$, independently of $J(s)$, so $\frac{\mathrm d}{\mathrm ds}\E J=\nu(A-I)\E J$ and $\E J(t)=e^{\nu t(A-I)}=e^{-(1-\bar J_{\mathrm{diag}})\nu t}e^{\nu tA_{\mathrm{off}}}$, whose entries, all nonnegative, are the series in \eqref{eq:sca-poisson}.
\end{proof}

\begin{proof}[Proof of \cref{prop:sca-routing}]
For $q(\mathbf g)_\pi\propto E_i(\pi)e^{-\gamma g_\pi}$ and $\mathbf g(s)=\mathbf g+s\boldsymbol\Delta$, $\frac{\mathrm d}{\mathrm ds}q_\pi=-\gamma q_\pi(\Delta_\pi-\E_q\Delta)$, so $\Vert\frac{\mathrm d}{\mathrm ds}q\Vert_1=\gamma\E_q|\Delta-\E_q\Delta|\le\gamma(\Var_q\Delta)^{1/2}\le\frac\gamma2\mathrm{osc}(\Delta)$ by Popoviciu's inequality $\Var\le\mathrm{osc}^2/4$; integrating over $s\in[0,1]$, $\Vert q(\mathbf g+\boldsymbol\Delta)-q(\mathbf g)\Vert_1\le\frac\gamma2\mathrm{osc}(\Delta)$. With $\Delta_\pi=G_i(\pi;\mu')-G_i(\pi;\mu)$ and $\delta=\mu'-\mu$, $\mathrm{osc}(\Delta)\le2\sum_k\mathrm{Lip}_{ik}\Vert\delta_k\Vert$.

(a) The averaged update changes by $\sum_\pi q'_\pi[\mathfrak U_{i,\pi}(\mu')-\mathfrak U_{i,\pi}(\mu)]+\sum_\pi(q'_\pi-q_\pi)\mathfrak U_{i,\pi}(\mu)$. The first sum is at most $\sum_k\bar A_{ik}\Vert\delta_k\Vert$ by the mean value inequality on the convex domain. The second has weights of total zero, so it is at most $\frac12\Vert q'-q\Vert_1\max_{\pi,\pi'}\Vert\mathfrak U_{i,\pi}(\mu)-\mathfrak U_{i,\pi'}(\mu)\Vert\le\frac\gamma2\mathrm{osc}_i\sum_k\mathrm{Lip}_{ik}\Vert\delta_k\Vert$. The proofs of \cref{thm:sca-lightcone}(c) and \cref{prop:sca-poisson} use only such one-step bounds, with differences in place of derivatives.

(b) Given coupled states, draw each agent's two policies from the maximal coupling of $q_i(\cdot\mid\mu)$ and $q_i(\cdot\mid\mu')$, independently across agents, and share all other noise. The policies differ with probability $\frac12\Vert q-q'\Vert_1\le\frac\gamma2\sum_k\mathrm{Lip}_{ik}\Vert\delta_k\Vert$; when they agree the update difference is at most $\sum_k\bar A_{ik}\Vert\delta_k\Vert$, and otherwise at most $\mathrm{osc}_i+\sum_k\bar A_{ik}\Vert\delta_k\Vert$. Hence $\E[\Vert\mu_i^+-\mu_i'^+\Vert\mid\mu,\mu']\le\sum_k\bar J_{ik}(\gamma)\Vert\delta_k\Vert$. Take expectations and iterate, using $\bar J(\gamma)\ge0$.

(c) A maximum of sums is at most the sum of the maxima, and \cref{prop:sca-poisson} with $\bar J_{\mathrm{diag}}(\gamma)\le1$ gives $v_{\mathrm{LR}}\le\inf_{\varphi>0}\nu g(\gamma)e^\varphi/\varphi=eg(\gamma)\nu$.

(d) By (a) the averaged map is a contraction with constant $g_{\mathrm{tot}}(\gamma)$ in $\max_i\Vert\cdot\Vert$, and the Banach fixed-point theorem gives a unique fixed point approached exponentially. By (b) the vector of expected coordinate distances contracts by $g_{\mathrm{tot}}(\gamma)$ per round in the maximum norm; on $N$ agents this bounds the Wasserstein distance between the laws at time $t$, for the cost $\sum_i\Vert\mu_i-\mu'_i\Vert$, by $Ng_{\mathrm{tot}}(\gamma)^t$ times the initial one, so some power of the transition operator is a strict contraction on the complete space of laws with finite first moment on the closed domain. The operator maps this space to itself: the one-step law from one state has a finite first moment by hypothesis, and by (b) with finite $\mathrm{osc}_i$ the one-step laws from two states lie at Wasserstein distance at most $\sum_{i,k}\bar J_{ik}(\gamma)\Vert\delta_k\Vert<\infty$, so the one-step law from every state, and from every initial law with finite first moment, has a finite first moment. This gives a unique stationary law reached at an exponential rate. The condition $g_{\mathrm{tot}}(\gamma)<1$ reads $\sum_k\bar A_{ik}+\frac\gamma2\mathrm{osc}_i\sum_k\mathrm{Lip}_{ik}<1$ for every $i$, that is $\gamma<\gamma_D$.
\end{proof}

\begin{proof}[Proof of \cref{thm:sca-no-moment}]
(a) Write $e(t)=\Phi(t,0)e(0)+\sum_{u<t}\Phi(t,u+1)\zeta(u)$; under \textup{(Lin)} the terms are independent and centered. Block $(i,j)$ of the term injected at $u$ is $\sum_k\Phi_{ik}Q^\zeta_k(u)\Phi_{jk}^\top$ with $n=t-1-u$ factors. By \cref{thm:sca-lightcone}(a) applied to the linear maps, it vanishes unless some $k$ has $d(i,k)\le n$ and $d(j,k)\le n$, which forces $n\ge\lceil d/2\rceil$. By \cref{thm:sca-lightcone}(c), $\Vert\Phi_{jk}\Vert\le(\bar J^n)_{jk}\le e^{-\varphi d(j,k)}g_\varphi^n$, and induction with $e^{\varphi d(i,k)}\le e^{\varphi d(i,l)}e^{\varphi d(l,k)}$ gives $\sum_k(\bar J^n)_{ik}e^{\varphi d(i,k)}\le g_\varphi^n$. With $e^{-\varphi d(j,k)}\le e^{-\varphi d(i,j)}e^{\varphi d(i,k)}$,
\begin{equation*}
\sum_k\Vert\Phi_{ik}\Vert\Vert\Phi_{jk}\Vert\le e^{-\varphi d(i,j)}g_\varphi^n\sum_k(\bar J^n)_{ik}e^{\varphi d(i,k)}\le e^{-\varphi d(i,j)}g_\varphi^{2n},
\end{equation*}
so the term is at most $qe^{-\varphi d}g_\varphi^{2n}$. Summing over $n$ from $\lceil d/2\rceil$ to $t-1$ and adding the initial term, an injection with $n=t$ and $q_0$ in place of $q$, gives \eqref{eq:sca-no-moment}. For the lag, $\Cov(e(t+h),e(t))=\Phi(t+h,t)\Cov(e(t))$ because later noise is independent of $e(t)$, so $\Cov(e_i(t+h),e_j(t))=\sum_k\Phi_{ik}(t+h,t)\Sigma_{kj}(t)$. Write $B(d')$ for the bracket of \eqref{eq:sca-no-moment} at distance $d'$; it is nonincreasing in $d'$, and $\Vert\Sigma_{kj}(t)\Vert\le e^{-\varphi d(k,j)}B(d(k,j))$. By \cref{thm:sca-lightcone}(a), $\Phi_{ik}(t+h,t)$ vanishes unless $d(i,k)\le h$, and then $d(k,j)\ge(d-h)_+$, so $B(d(k,j))\le B((d-h)_+)$. The same inequality as above gives $\sum_k\Vert\Phi_{ik}(t+h,t)\Vert e^{-\varphi d(k,j)}\le e^{-\varphi d(i,j)}g_\varphi^h$, hence the lagged covariance is at most $e^{-\varphi d}g_\varphi^hB((d-h)_+)$, which is \eqref{eq:sca-no-moment-lag} because $\one\{2t\ge(d-h)_+\}=\one\{2t+h\ge d\}$.

(b) $\Sigma_{ii}(t)\succeq Q^\zeta_i(t-1)\succeq q_{\min}I$ because the last injection is independent of everything else, so the top canonical correlation $\Vert\Sigma_{ii}^{-1/2}\Sigma_{ij}\Sigma_{jj}^{-1/2}\Vert$ is at most $\Vert\Sigma_{ij}\Vert/q_{\min}$. For $g_\varphi>1$, $\sum_{n<t}g_\varphi^{2n}\le g_\varphi^{2t}/(g_\varphi^2-1)$, so the correlation is at most $K_\varphi e^{-\varphi d+2t\ln g_\varphi}=K_\varphi\exp(-\varphi[d-2v_\varphi t])$, which is below $\theta$ once $d>2v_\varphi t+\varphi^{-1}\ln(K_\varphi/\theta)$. The bound $2t$ is the support statement.

(c) For $g_\varphi<1$, $\sum_{n\ge\lceil d/2\rceil}g_\varphi^{2n}\le g_\varphi^d/(1-g_\varphi^2)$, and $g_\varphi^{2t}\le g_\varphi^d$ whenever $2t\ge d$, so the bound is at most $e^{-(\varphi+\ln(1/g_\varphi))d}(q/(1-g_\varphi^2)+q_0)$ for every $t$.
\end{proof}

\begin{proof}[Proof of \cref{prop:sca-rank-one}]
Let $r=\max_{k\ge2}|\nu_k|<\rho$ over the other eigenvalues $\nu_k$ of $W$. Then $W^n=\rho^nvu^\top+R_n$ with $R_n=W^n(I-vu^\top)$ and $\Vert R_n\Vert\le C_\delta(r+\delta)^n$ for every $\delta>0$. In $\Sigma(t)=W^t\Sigma(0)W^{t\top}+\sum_{n<t}W^nQ^\zeta W^{n\top}$, the leading terms sum to $\rho^{2t}[u^\top\Sigma(0)u+u^\top Q^\zeta u(1-\rho^{-2t})/(\rho^2-1)]vv^\top$. The cross terms are $O(t\max\{\rho(r+\delta),1\}^t)$ and the remainders $O(t\max\{(r+\delta)^2,1\}^t)$; with $\delta<\rho-r$ both are $o(\rho^{2t})$. This proves \eqref{eq:sca-rank-one}, and $c_\infty>0$ because $Q^\zeta\succ0$. For the readouts $\hat v_i^\top e_i$ with $\hat v_i=v_i/|v_i|$, covariances and variances are $\rho^{2t}c_\infty|v_i||v_j|(1+o(1))$ and $\rho^{2t}c_\infty|v_i|^2(1+o(1))$, so their correlation tends to one, and so does the top canonical correlation. The top eigenvalue of $\Sigma(t)$ is $c_\infty\rho^{2t}|v|^2(1+o(1))$ with eigenvector tending to $v/|v|$, which gives the eigenvalue of order $\rho^{-2t}$ of the precision. For the others, the Courant--Fischer theorem bounds $\lambda_2(\Sigma(t))$ by the largest Rayleigh quotient on $v^\perp$. There the leading terms vanish, and so do the cross terms, which have the form $\rho^n(va_n^\top+a_nv^\top)$ with $a_n=R_nQ^\zeta u$ and its analog for $\Sigma(0)$, since $x^\top v=0$. What remains is bounded by $\sum_n\Vert R_n\Vert^2\Vert Q^\zeta\Vert+\Vert R_t\Vert^2\Vert\Sigma(0)\Vert$, which is bounded in $t$ when $r<1$. So $\lambda_2(\Sigma(t))$ is bounded, and every eigenvalue of $\Sigma(t)^{-1}$ but the smallest is bounded away from zero.
\end{proof}

\begin{proof}[Proof of \cref{prop:sca-amplifying-site}]
(a) Away from the origin the eigenvalue equation is $\rho v_x=cv_x+a(v_{x+1}+v_{x-1})$, whose square-summable solutions are multiples of $e^{-m|x|}$ with $c+2a\cosh m=\rho$; at the origin, $\rho=c+b_0+2ae^{-m}$. Hence $2a\sinh m=b_0$ and $\rho=c+2a\cosh(\arcsinh(b_0/(2a)))=c+\sqrt{4a^2+b_0^2}$. The operator $W$ is self-adjoint, $W\succeq(c-2a)I$ since $b_0>0$, and every two-dimensional subspace contains a unit vector with vanishing component at the origin, on which the Rayleigh quotient is that of the medium, at most $c+2a$. By the min-max principle at most one eigenvalue exceeds $c+2a$, and the spectrum of $W$ on $v^\perp$ lies in $[c-2a,c+2a]$.

(b) Write $W=\rho P_v+W_c$ with $W_c=(I-P_v)W(I-P_v)$, which commutes with $P_v$ and has norm at most $\max(|c-2a|,c+2a)=\lambda$, since $c\ge0$. Then $W^{2n}=\rho^{2n}P_v+W_c^{2n}$, with the convention $W_c^0=I-P_v$, and from $\Sigma(0)=0$, $\Sigma(t)=q\sum_{n<t}W^{2n}=A(t)P_v+B(t)$ with $B(t)=q\sum_{n<t}W_c^{2n}$, so $q(I-P_v)\preceq B(t)\preceq q(1-\lambda^2)^{-1}(I-P_v)$. Also $|v|^2=\sum_xe^{-2m_\rho|x|}=\coth m_\rho$.

(c) Since $W\ge0$ entrywise, $\Sigma(t)\ge0$ entrywise, so the lower bound holds trivially when $\mathfrak s_x\le1$. Put $\beta=q/(1-\lambda^2)$, $a_1=Av_x^2/|v|^2$ and $a_0=A/|v|^2$, where $v_0=1\ge v_x$. Since $|B_{xy}|\le\beta$, $\Sigma_{x0}\ge\sqrt{a_1a_0}-\beta$, $\Sigma_{xx}\le a_1+\beta$ and $\Sigma_{00}\le a_0+\beta$. For $\mathfrak s_x>1$ the numerator is positive, the function $a_0\mapsto(\sqrt{a_1a_0}-\beta)/\sqrt{(a_1+\beta)(a_0+\beta)}$ is increasing (its derivative in $\sqrt{a_0}$ has numerator $\beta(\sqrt{a_1}+\sqrt{a_0})>0$), and at $a_0=a_1$ it equals $(a_1-\beta)/(a_1+\beta)=(\mathfrak s_x-1)/(\mathfrak s_x+1)$, which gives the lower bound. For the upper bound use $\Sigma_{xx}\ge q$, from the fresh noise, and $\Sigma_{00}\ge a_0$, from $B(t)\succeq0$: $\Corr\le(\sqrt{a_1a_0}+\beta)/\sqrt{qa_0}$, which is the right side of \eqref{eq:sca-front-bounds}.

(d) For $x\ge0$, $\mathfrak s_x(t)=(1-\lambda^2)(\rho^{2t}-1)e^{-2m_\rho x}/((\rho^2-1)\coth m_\rho)$. The lower bound gives correlation at least $\theta$ when $\mathfrak s_x\ge(1+\theta)/(1-\theta)$, which, since $\rho^{2t}-1\ge(1-\rho^{-2})\rho^{2t}$ for $t\ge1$, holds for $x\le(t\ln\rho-c_1(\theta))/m_\rho$ with a constant $c_1(\theta)$. The upper bound gives correlation below $\theta$ when both of its terms are below $\theta/2$; the second is for $t\ge t_0(\theta)$ because $A(t)\to\infty$, and the first is when $\mathfrak s_x<\theta^2(1-\lambda^2)/4$, which holds for $x>(t\ln\rho+c_2(\theta))/m_\rho$ with a constant $c_2(\theta)$. Hence $x_\theta(t)=t\ln\rho/m_\rho+O_\theta(1)$. Finally $\rho=c+2a\cosh m_\rho\le(c+2a)\cosh m_\rho<\cosh m_\rho<e^{m_\rho}$, so $\ln\rho<m_\rho$.

(e) Outside the core the eigenvector solves the medium's recurrence $\rho v_x=cv_x+a(v_{x+1}+v_{x-1})$, whose solutions are combinations of $e^{\pm m_\rho x}$ with $c+2a\cosh m_\rho=\rho$; square summability selects the decaying one on each side. As $\rho\downarrow1$, $(\rho-c)/(2a)$ decreases to $(1-c)/(2a)=1+\varepsilon/(2a)$, so $m_\rho$ decreases to the rate of \cref{thm:sca-oz}(e).
\end{proof}

\begin{proof}[Proof of \cref{thm:sca-clustering}]
(a) The support statement is \cref{thm:sca-lightcone}(a) for linear maps, and the row-sum bound is \cref{lem:sca-relations}(a). (b) \textup{(M)} is \cref{lem:sca-relations}(a). Block $(i,j)$ of the $n$-th term of the series is $\sum_k\Phi_{ik}Q^\zeta_k\Phi_{jk}^\top$ with $\Phi=\Phi(t,t-n)$. It vanishes unless $d(i,k)\le n$ and $d(j,k)\le n$ for some $k$, which forces $n\ge\lceil d(i,j)/2\rceil$, and its norm is at most $q\sum_k\Vert\Phi_{ik}\Vert\Vert\Phi_{jk}\Vert\le q(\max_k\Vert\Phi_{jk}\Vert)\sum_k\Vert\Phi_{ik}\Vert\le q\lambda^{2n}$. Summing the geometric series from $\lceil d/2\rceil$ gives \eqref{eq:sca-clustering}, with $\lambda^{2\lceil d/2\rceil}\le\lambda^d=e^{-m_\Sigma d}$, and absolute convergence. Every term of the series is positive semidefinite and the term $n=0$ is $Q^\zeta(t-1)$, so $\Sigma(t)\succeq Q^\zeta(t-1)$; under \textup{(B)} this gives $\lambda_{\min}(\Sigma(t))\ge q_{\min}$, $\Sigma_{ii}\succeq q_{\min}I$, hence $\Vert R_{ij}\Vert\le\Vert\Sigma_{ij}\Vert/q_{\min}$ and $\xi\le1/\ln(1/\lambda)$. (c) The initial covariance is one more injection, $n=t-s$ steps back, with $q_0$ in place of $q$. (d) Later noise is independent of $e(t)$, so $\Cov(e(t+h),e(t))=\Phi(t+h,t)\Sigma(t)$; block $(i,j)$ is at most $\sum_k\Vert\Phi_{ik}\Vert\Vert\Sigma_{kj}\Vert$ over $k$ with $d(i,k)\le h$, hence $d(k,j)\ge d(i,j)-h$, and the sum is at most $\lambda^h\cdot q\lambda^{(d-h)_+}/(1-\lambda^2)=q\lambda^{\max(h,d)}/(1-\lambda^2)$.
\end{proof}

\begin{proof}[Proof of \cref{prop:sca-world-conditioning}]
(a) Block $(i,j)$ of the $n$-th term is $\sum_{k,l}\Phi_{ik}Q^\zeta_{kl}\Phi_{jl}^\top$, with $d(i,k)\le n$ and $d(j,l)\le n$ on the support, so $d(k,l)\ge d-2n$ and $\Vert Q^\zeta_{kl}\Vert\le qe^{-(d-2n)_+/\xi_w}$. Since $\sum_{k,l}\Vert\Phi_{ik}\Vert\Vert\Phi_{jl}\Vert\le\lambda^{2n}$, the term is at most $q\lambda^{2n}e^{-(d-2n)_+/\xi_w}$. Each of the at most $\lceil d/2\rceil$ terms with $2n<d$ is at most $q\max(\lambda,e^{-1/\xi_w})^d$, and the terms with $2n\ge d$ sum to at most $q\lambda^d/(1-\lambda^2)$. (b) Without messages $W=\lambda I$ with $\lambda=F(1-k^s)$, so the connected part has the diagonal covariance $qI/(1-\lambda^2)$, and by \cref{lem:fw-common-mode}(d) the common mode adds $Q\one\one^\top/(1-\lambda^2)$, independent of the connected part. Every error then has variance $(Q+q)/(1-\lambda^2)$, and any two distinct errors have covariance $Q/(1-\lambda^2)$.
\end{proof}

\begin{proof}[Proof of \cref{prop:sca-spectral}]
The stationary covariance is $\sum_nW^nQ^\zeta W^{n\top}$. Block $(i,j)$ of the $n$-th term vanishes for $n<\lceil d(i,j)/2\rceil$, as in the proof of \cref{thm:sca-clustering}, and its norm is at most $\Vert W^nQ^\zeta W^{n\top}\Vert\le q\Vert W^n\Vert^2\le qK_A^2f^{2n}$; sum the series. For normal $W$, $\Vert W^n\Vert=\rho(W)^n$. For every $f>\rho(W)$ there is $K_A$ with $\Vert W^n\Vert\le K_Af^n$ (Gelfand's formula), so with $\Vert R_{ij}\Vert\le\Vert\Sigma_{ij}\Vert/q_{\min}$ the decay rate of correlations is at least $\ln(1/f)$ for every $f>\rho(W)$, which gives $\xi\le\tau$.
\end{proof}

\begin{proof}[Proof of \cref{thm:sca-diffusive}]
We first prove Carne's bound. For $w\in\R$ and $z=\cos w$,
\begin{equation*}
z^n=\E\cos(S_nw)=\sum_{k=-n}^n\Pr(S_n=k)T_{|k|}(z),
\end{equation*}
where $S_n$ is the simple random walk on $\Z$ and $T_k$ is the Chebyshev polynomial with $T_k(\cos w)=\cos(kw)$. Reversibility makes $P$ self-adjoint on $L^2(\varpi)$ with spectrum in $[-1,1]$, so $P^n=\sum_k\Pr(S_n=k)T_{|k|}(P)$ with $\Vert T_k(P)\Vert_{L^2(\varpi)}\le1$, and $\langle\one_x,T_k(P)\one_y\rangle_\varpi=0$ for $k<d(x,y)$ because $T_k$ has degree $k$ and $(P^l)_{xy}=0$ for $l<d(x,y)$. Hence $\varpi_xP^n_{xy}=\langle\one_x,P^n\one_y\rangle_\varpi\le\Pr(|S_n|\ge d(x,y))\sqrt{\varpi_x\varpi_y}$. By Markov's inequality applied to $e^{sS_n}$, with $\E e^{sS_n}=\cosh^ns\le e^{ns^2/2}$ and $s=d/n$, $\Pr(|S_n|\ge d)\le2e^{-d^2/(2n)}$, so $P^n_{xy}\le2\sqrt{\varpi_y/\varpi_x}e^{-d(x,y)^2/(2n)}$.

Clustering. By induction, $\Vert\Phi_{ik}(t,t-n)\Vert\le\lambda^n(P^n)_{ik}$, so the series of \cref{thm:sca-clustering}(b) for the process started in the infinite past gives $\Vert\Sigma_{ij}(t)\Vert\le q\sum_n\lambda^{2n}\sum_kP^n_{ik}P^n_{jk}$; from a zero start the sum over $n$ is finite and the same bound holds. Reversibility gives $P^n_{jk}=\varpi_kP^n_{kj}/\varpi_j$, so $\sum_kP^n_{ik}P^n_{jk}\le(\varpi_{\max}/\varpi_j)P^{2n}_{ij}\le2\varpi_{\max}(\varpi_i\varpi_j)^{-1/2}e^{-d^2/(4n)}$. With $\varepsilon_\lambda=\ln(1/\lambda)$, the exponent $h(n)=2n\varepsilon_\lambda+d^2/(4n)$ satisfies $h(n)\ge d\sqrt{2\varepsilon_\lambda}$ by the inequality between arithmetic and geometric means, hence $h(n)\ge\vartheta d\sqrt{2\varepsilon_\lambda}+(1-\vartheta)2n\varepsilon_\lambda$, and the geometric sum over $n$ gives \eqref{eq:sca-diffusive}.
\end{proof}

\begin{proof}[Proof of \cref{prop:sca-heterogeneous}]
(a) Let $(W_{\boldsymbol\varphi})_{ij}=|W_{ij}|e^{\boldsymbol\varphi\cdot(x_i-x_j)}$. The moduli of the entries of a product are at most the entries of the product of the moduli, and the tilts telescope along products, so the tilted moduli $|\Phi_{ik}(t,t-n)|e^{\boldsymbol\varphi\cdot(x_i-x_k)}$ are at most the entries of a product of $n$ matrices $W_{\boldsymbol\varphi}(s)$, whose row sums are at most $\bar\psi(\boldsymbol\varphi)^n$. For scalar agents,
\begin{equation*}
|\Sigma_{ij}|e^{\boldsymbol\varphi\cdot(x_i-x_j)}\le q_{\max}\sum_n\sum_k|\Phi_{ik}|e^{\boldsymbol\varphi\cdot(x_i-x_k)}|\Phi_{jk}|e^{-\boldsymbol\varphi\cdot(x_j-x_k)}\le q_{\max}\sum_n\bar\psi(\boldsymbol\varphi)^n\bar\psi(-\boldsymbol\varphi)^n,
\end{equation*}
using that the largest entry of a nonnegative row is at most its sum; the same holds for any finite sum, hence from a zero start. The telescoping uses single-valued positions, which $\Z^D$ and a box with open ends have, and the row sums are taken over all agents, boundary agents included. On a periodic box of side $L$ let $\tilde W(t)$ be the $L$-periodic medium on $\Z^D$ with $\tilde W_{\tilde\imath\tilde\jmath}=W_{ij}$ for neighboring lifts $\tilde\imath,\tilde\jmath$ of $i,j$, driven by independent noises with the periodically extended variances, and let $\tilde\Phi$ and $\tilde\Sigma$ be its propagator and covariance. Every walk on the box from $i$ to $k$ lifts uniquely to a walk from a fixed lift $\tilde\imath$ to one lift of $k$, so $\Phi_{ik}=\sum_{\mathbf m\in\Z^D}\tilde\Phi_{\tilde\imath,\tilde k+L\mathbf m}$, and by periodicity $\tilde\Phi_{\tilde\jmath,\tilde k+L\mathbf m'}=\tilde\Phi_{\tilde\jmath+L(\mathbf m-\mathbf m'),\tilde k+L\mathbf m}$. Hence
\begin{equation*}
\Sigma_{ij}=\sum_n\sum_kQ^\zeta_k\Phi_{ik}\Phi_{jk}=\sum_{\mathbf w\in\Z^D}\sum_n\sum_{\tilde k}Q^\zeta_{\tilde k}\tilde\Phi_{\tilde\imath\tilde k}\tilde\Phi_{\tilde\jmath+L\mathbf w,\tilde k}=\sum_{\mathbf w\in\Z^D}\tilde\Sigma_{\tilde\imath,\tilde\jmath+L\mathbf w},
\end{equation*}
with $\tilde\jmath+L\mathbf w$ the lift of $j$ displaced by $L\mathbf w$. The rearrangement is justified by absolute convergence: $\sum_{\mathbf w}|\tilde\Phi_{\tilde\jmath+L\mathbf w,\tilde k}|=\sum_{\mathbf w}|\tilde\Phi_{\tilde\jmath,\tilde k-L\mathbf w}|$ is part of a row sum of the product of the moduli, hence at most $\bar\psi(0)^n$, and $\bar\psi(0)^2\le\Psi(\boldsymbol\varphi)<1$ by the Cauchy--Schwarz inequality applied to each row. Each term obeys the bound for $\tilde W$ with any admissible $\boldsymbol\varphi$. On $\Z^D$, for $W(t)\ge0$ with equal weights $a_{id}$ to the two neighbors along axis $d$, $\sum_jW_{ij}e^{\boldsymbol\varphi\cdot(x_i-x_j)}=W_{ii}+\sum_d2a_{id}\cosh\varphi_d=1-\varepsilon_i+2\sum_da_{id}(\cosh\varphi_d-1)$, which is at most $1-\varepsilon_{\min}+2a_{\max}\sum_d(\cosh\varphi_d-1)$ and is even in $\boldsymbol\varphi$. Along an axis, $\boldsymbol\varphi=m\mathbf u_d$ gives $\Psi<1$ when $\cosh m<1+\varepsilon_{\min}/(2a_{\max})$; along the diagonal, $\varphi_d=s$ for all $d$ gives $\boldsymbol\varphi\cdot\mathbf u=\sqrt Ds$ and $\Psi<1$ when $\cosh s<1+\varepsilon_{\min}/(2Da_{\max})$.

(b) Products of nonnegative matrices are monotone in each factor, so $\Phi(t,t-n)\ge W_{\mathrm{low}}^n$ entrywise, and with the diagonal $Q^\zeta\ge q_{\min}I$, $\Sigma\ge q_{\min}\sum_nW_{\mathrm{low}}^{2n}=q_{\min}(I-W_{\mathrm{low}}^2)^{-1}$. Each row of $W(t)$ dominates the corresponding row of $W_{\mathrm{low}}$, so $1-\varepsilon_i\ge c_{\min}+2Da_{\min}=1-\varepsilon_{\mathrm{low}}$ for an agent with $2D$ neighbors. A boundary agent of an open box with $n_i<2D$ neighbors only has $1-\varepsilon_i\ge c_{\min}+n_ia_{\min}$; for $W=W_{\mathrm{low}}$ on an open chain the end leak is $1-c_{\min}-a_{\min}>\varepsilon_{\mathrm{low}}$.

(c) On $\Z$ the hypotheses of the second part of (a) and of (b) hold. By (a) at $\boldsymbol\varphi=0$, $\Sigma_{ii}\le q_{\max}/(1-\bar\psi(0)^2)$, and $\Sigma_{ii}\ge q_{\min}$. So $R_{ij}\le|\Sigma_{ij}|/q_{\min}$ decays at least at the rate $m_{\mathrm{up}}$, and $R_{ij}\ge q_{\min}(1-\bar\psi(0)^2)(I-W_{\mathrm{low}}^2)^{-1}_{ij}/q_{\max}$ decays at most at the rate $m_{\mathrm{low}}$ of the homogeneous medium, which in one dimension is exact by \cref{thm:sca-oz}(e).

(d) For a nonnegative matrix the spectral radius lies between the smallest and the largest row sum (Perron--Frobenius).

(e) The definitions give $W_{ij}\le\lambda'P_{ij}$ for all $i$, $j$, with $P$ symmetric and stochastic and $P_{ii}=1-\sum_{j\neq i}P_{ij}\ge W_{ii}/\lambda'\ge0$. A symmetric chain is reversible with respect to the uniform measure, which is the hypothesis of \cref{thm:sca-diffusive}.
\end{proof}

\begin{proof}[Proof of \cref{thm:sca-oz}]
(a) $W$ is symmetric with $\Vert W\Vert=\sup_p|\hat W(p)|=\lambda<1$, so $\Sigma=q\sum_nW^{2n}=q(I-W^2)^{-1}$, and $(I-W^2)^{-1}=\frac12[(I-W)^{-1}+(I+W)^{-1}]$.

(b) Since $c,a\ge0$ the kernel is nonnegative and $\sum_xW^n(x)e^{\boldsymbol\varphi\cdot x}=\psi(\boldsymbol\varphi)^n$. By Tonelli's theorem $\sum_x\Sigma(x)e^{\boldsymbol\varphi\cdot x}=q\sum_n\psi(\boldsymbol\varphi)^{2n}=q/(1-\psi(\boldsymbol\varphi)^2)$, and each nonnegative term is at most the sum. For $\boldsymbol\varphi$ in the interior of $\{\psi\le1\}$, $-\ln\Sigma(x)\ge\boldsymbol\varphi\cdot x-\ln(q/(1-\psi^2))$; take $\boldsymbol\varphi$ close to the maximizer of $\boldsymbol\varphi\cdot\mathbf u$.

(c) Summing a convolution power over the transverse coordinates gives the convolution power of the summed kernel, in which transverse moves become stays: $\sum_yW^n(t,y)=W_1^n(t)$, where $W_1$ is the one-dimensional kernel with self-weight $c'=c+2(D-1)a$ and neighbor weight $a$, with the same $\lambda$. So the hyperplane sum is the one-dimensional covariance of (e) with $c'$, which decays at the rate $m_1$ with $\cosh m_1=(1-c')/(2a)=1+\varepsilon/(2a)$; its staggered part decays faster when $c'>0$ and vanishes on odd $t$ when $c'=0$, so in all cases $-\limsup t^{-1}\ln\sum_y\Sigma(t,y)=m_1$. For the support function along $\mathbf u_1$, the constraint $c+2a\sum_d\cosh\varphi_d\le1$ is least restrictive for $\varphi_1$ when $\varphi_d=0$ for $d\ge2$, which gives $\cosh\varphi_1\le1+\varepsilon/(2a)$. Since $\cosh m-1=2\sinh^2(m/2)$, $m_1=2\arcsinh\sqrt{\varepsilon/(4a)}$, and $\arccosh(1+y)=\sqrt{2y}(1-y/12+O(y^2))$ with $y=\varepsilon/(2a)$ gives the expansion.

(d) The set $\{\psi\le1\}$ is convex and invariant under permutations of the coordinates, and so is the objective $\boldsymbol\varphi\cdot\mathbf u$ for the diagonal $\mathbf u$; averaging a maximizer over permutations gives a symmetric maximizer, $\varphi_d=s$ with $c+2aD\cosh s=1$, that is $\cosh s=1+\varepsilon/(2Da)$, and $\boldsymbol\varphi\cdot\mathbf u=\sqrt Ds$. The expansion of (c) with $y=\varepsilon/(2Da)$ gives $m(\mathbf u)=\sqrt{\varepsilon/a}(1-\varepsilon/(24Da)+O(\varepsilon^2))$, and the anisotropy follows by subtraction.

(e) In one dimension $1-\hat W(p)=2a(\cosh m-\cos p)$ and $1+\hat W(p)=2a(\cosh m'+\cos p)$, with $\cosh m=(1-c)/(2a)=1+\varepsilon/(2a)>1$ and $\cosh m'=(1+c)/(2a)>1$ because $2a<1-c\le1+c$. The residue computation $\int_{-\pi}^{\pi}e^{ipx}(\cosh m-\cos p)^{-1}\mathrm{d}p/2\pi=e^{-m|x|}/\sinh m$, and the shift $p\to p+\pi$ for the second resolvent, give \eqref{eq:sca-oz-1d} through (a). On a ring of even length $N$ the covariance is the periodization $\sum_{j\in\Z}\Sigma(x+jN)$, and $\sum_je^{-m|x+jN|}=\cosh(m(N/2-|x|))/\sinh(mN/2)$ for $|x|\le N/2$, with $(-1)^{x+jN}=(-1)^x$.

(f) $\Sigma(0)=\frac q2\int[(1-\hat W)^{-1}+(1+\hat W)^{-1}]\mathrm{d}p/(2\pi)^D$. With $c>0$, $1+\hat W\ge1+c-2Da=2c+\varepsilon$, so the second integral is bounded. Write $1-\hat W(p)=\varepsilon+ag(p)$ with $g(p)=2\sum_d(1-\cos p_d)=|p|^2(1+O(|p|^2))$, and $g$ bounded below away from the origin. For $D=1$ the claim follows from (e), since $\sinh m\sim\sqrt{\varepsilon/a}$ while $m'$ stays bounded away from zero. For $D=2$, the difference between $\int_{|p|\le1}(\varepsilon+ag)^{-1}$ and $\int_{|p|\le1}(\varepsilon+a|p|^2)^{-1}=(\pi/a)\ln((\varepsilon+a)/\varepsilon)$ is bounded uniformly in $\varepsilon$, as is the integral over $|p|>1$; dividing by $(2\pi)^2$ and multiplying by $q/2$ gives $(q/8\pi a)\ln(1/\varepsilon)+O(1)$. For $D\ge3$, $(ag)^{-1}$ is integrable and monotone convergence gives a finite limit.

(g) By (b), $\Sigma$ decays exponentially, so $\hat\Sigma$ is smooth and $\sum_x|x|^2\Sigma(x)=-\Delta_p\hat\Sigma(0)$. From $\hat W(p)=\lambda-a|p|^2+O(|p|^4)$, $\hat\Sigma(p)=q/(1-\lambda^2)-2aq\lambda|p|^2/(1-\lambda^2)^2+O(|p|^4)$, so $-\Delta_p\hat\Sigma(0)=4Dqa\lambda/(1-\lambda^2)^2$, and $\sum_x\Sigma(x)=\hat\Sigma(0)=q/(1-\lambda^2)$.
\end{proof}

\begin{proof}[Proof of \cref{prop:sca-lengths-relation}]
(a) For a translation-invariant scalar field, $\hat\Sigma=1/\hat\Pi$. By symmetry $\partial_d\hat\Pi(0)=0$, and $\partial_d^2\hat\Pi(0)=\varsigma_d^2$, $\hat\Pi(0)=\varepsilon^\Pi$, so $\partial_d^2\hat\Sigma(0)=-\varsigma_d^2/(\varepsilon^\Pi)^2$. Hence $\sum_xx_d^2\Sigma(x)=\varsigma_d^2/(\varepsilon^\Pi)^2$ and $\sum_x\Sigma(x)=1/\varepsilon^\Pi$, and the ratio is $2\xi_{2,d}^2$. The walk expansion $\Sigma=\sum_m\mathcal A^m/(\varepsilon^\Pi+A)^{m+1}$, with $\mathcal A$ the convolution by $a^\Pi$ and $A=\sum_xa^\Pi(x)$, shows that $\Sigma\ge0$ has finite second moments, which justifies the differentiation.

(b) By Tonelli's theorem and the walk expansion, $\sum_xe^{\vartheta x_d}\Sigma(x)=(\varepsilon^\Pi-\sum_xa^\Pi(x)(\cosh(\vartheta x_d)-1))^{-1}$ when the bracket is positive and $+\infty$ otherwise. The left side is a two-sided series in $e^\vartheta$ with nonnegative coefficients $g(t)=\sum_{x:x_d=t}\Sigma(x)=g(-t)$, finite for $|\vartheta|<m_d$ and tending to $+\infty$ as $\vartheta\uparrow m_d$, where $m_d$ is the positive root (it exists because the left side of the root equation is continuous, increasing, zero at zero and unbounded for couplings of finite range with $\varsigma_d>0$). So the power series $\sum_{t\ge0}g(t)z^t$ converges for $0<z<e^{m_d}$ and diverges at $z=e^{m_d}$ by monotone convergence; its radius of convergence is $e^{m_d}$, and the Cauchy--Hadamard formula gives $\limsup g(t)^{1/t}=e^{-m_d}$, the rate of \cref{def:sca-lengths}(b). Comparing Taylor coefficients, $y^2/2\le\cosh y-1\le(y^2/2)\cosh y$, so at $m=m_d$, $m^2\varsigma_d^2/2\le\varepsilon^\Pi\le(m^2\varsigma_d^2/2)\cosh(\ell^\Pi_dm)$. The first inequality is $m_d\le1/\xi_{2,d}$, and the second, with $\cosh(\ell^\Pi_dm_d)\le\cosh(\ell^\Pi_d/\xi_{2,d})$, gives $m_d\ge\xi_{2,d}^{-1}\cosh(\ell^\Pi_d/\xi_{2,d})^{-1/2}$.

(c) For nearest-neighbor couplings $\varsigma_d^2=2a^\Pi$, and the equation of (b) reads $2a^\Pi(\cosh m-1)=\varepsilon^\Pi=4a^\Pi\sinh^2(m/2)$, so $\xi_{2,d}^2=a^\Pi/\varepsilon^\Pi=(2\sinh(m/2))^{-2}$; the expansion $2\sinh(m/2)=m(1+m^2/24+O(m^4))$ gives the difference. For the medium of \cref{thm:sca-oz}, the world-conditional precision is $(I-W^2)/q$ with couplings $2ca/q$ at axis distance one, $a^2/q$ at axis distance two and $2a^2/q$ at the diagonal neighbors (\cref{thm:col-diffusive-precision}), so $\varsigma_d^2=(4ca+8a^2+8(D-1)a^2)/q=4a\lambda/q$ and $\varepsilon^\Pi=(1-\lambda^2)/q$, which gives $\xi_2^2=2a\lambda/(1-\lambda^2)$, in agreement with \cref{thm:sca-oz}(g). The hyperplane rate is $m$ with $\cosh m=1+\varepsilon/(2a)$ by \cref{thm:sca-oz}(c), and $\varepsilon=2a(\cosh m-1)=4a\sinh^2(m/2)$, so $a/\varepsilon=(2\sinh(m/2))^{-2}$ and $\xi_2^2=(2\lambda/(1+\lambda))(a/\varepsilon)$. The inequalities follow from $2\lambda/(1+\lambda)\le1$ and $2\sinh(m/2)\ge m$.
\end{proof}

\begin{proof}[Proof of \cref{prop:sca-persistence}]
$\Cov(e(t+h),e(t))=W^h\Sigma$, so with $u=\Sigma^{1/2}\mathbf v$ and $u'=\Sigma^{1/2}\mathbf v'$ the correlation is $u'^\top\tilde W^hu/(|u'||u|)$, whose supremum over $u$ and $u'$ is $\Vert\tilde W^h\Vert$; this supremum is attained, so $T^\ast(\theta)=\max\{h:\Vert\tilde W^h\Vert\ge\theta\}$. Also $\tilde W\tilde W^\top=\Sigma^{-1/2}W\Sigma W^\top\Sigma^{-1/2}=I-\Sigma^{-1/2}Q^\zeta\Sigma^{-1/2}$, which gives $\Vert\tilde W\Vert^2$. $\tilde W$ is similar to $W$, so $\rho(\tilde W)=\rho(W)$, and Gelfand's formula gives $\Vert\tilde W^h\Vert^{1/h}\to\rho(W)$. Hence $\rho(W)^h\le\Vert\tilde W^h\Vert\le C_\delta(\rho(W)+\delta)^h$ for every $\delta>0$, so $T^\ast(\theta)\ge\ln(1/\theta)/\ln(1/\rho(W))-1$ and $T^\ast(\theta)\le(\ln(1/\theta)+\ln C_\delta)/\ln(1/(\rho(W)+\delta))$. Dividing by $\ln(1/\theta)$ and letting $\theta\to0$ and then $\delta\to0$ gives the limit $\tau$; if $\rho(W)=0$, $\tilde W$ is nilpotent, $T^\ast$ stays bounded and the limit is $0=\tau$.
\end{proof}

\begin{proof}[Proof of \cref{prop:sca-diffusive-scaling}]
(a) $|\hat W(p)|\le c+2Da=\lambda$ with equality at $p=0$, so the spectral radius of the convolution is $\lambda$. With \cref{thm:sca-oz}(g), $\xi_2^2/\tau=2a\lambda\ln(1/\lambda)/(1-\lambda^2)$, and $\ln(1/\lambda)=\varepsilon+O(\varepsilon^2)$, $1-\lambda^2=2\varepsilon+O(\varepsilon^2)$. (b) $\Cov(e(t+h),e(t))=W^h\Sigma$ gives \eqref{eq:sca-autocov-def}. With $c>0$, $|\hat W(p)|=\lambda$ only at $p=0$, because $\hat W(p)=-\lambda$ would need $c+2a\sum_d\cos p_d=-c-2Da$, impossible for $c>0$; by continuity and compactness $|\hat W|\le\lambda-\eta_0$ outside any neighborhood of the origin, for some $\eta_0>0$, and that region contributes $O((\lambda-\eta_0)^h)$. Near the origin $\hat W(p)^h=\lambda^he^{-ha|p|^2/\lambda}(1+O(h|p|^4))$ and $(1-\hat W^2)^{-1}=(1-\lambda^2)^{-1}+O(|p|^2)$, and Laplace's method gives $\int e^{-ha|p|^2/\lambda}\mathrm{d}p/(2\pi)^D=(\lambda/4\pi ah)^{D/2}$ with relative error $O(1/h)$. The $O(|p|^2)$ term of $(1-\hat W^2)^{-1}$ is $-2\lambda a|p|^2/(1-\lambda^2)^2$, and the Gaussian weight has $\langle|p|^2\rangle=D\lambda/(2ah)$, so this term contributes the relative correction $-D\lambda^2/((1-\lambda^2)h)$, about $-D/(2\varepsilon h)$ for small $\varepsilon$; the remaining $O(1/h)$ corrections have coefficients bounded as $\varepsilon\to0$ at fixed $c$, so \eqref{eq:sca-autocov} is accurate once $\varepsilon h\gg1$.
\end{proof}

\begin{proof}[Proof of \cref{prop:sca-records}]
(a) With $\mathcal T=\diag(1,(\beta_1/\omega_1)^{1/2},\dots)$, $\mathcal T\mathsf R\mathcal T^{-1}$ is a symmetric arrowhead matrix with off-diagonal entries $-(\beta_k\omega_k)^{1/2}$, so the spectrum is real. (b) For $x\notin\{\omega_k+\delta_k\}$, the Schur complement of the diagonal record block gives
\begin{equation*}
\det(\mathsf R-xI)=\prod_k(\omega_k+\delta_k-x)\Bigl[\varepsilon'+\sum_k\beta_k-x-\sum_k\frac{\beta_k\omega_k}{\omega_k+\delta_k-x}\Bigr]=g_z(x)\prod_k(\omega_k+\delta_k-x),
\end{equation*}
so the eigenvalues below $d_{\min}=\min_k(\omega_k+\delta_k)$ are the roots of $g_z$ there. Since $g_z'(x)=-1-\sum_k\beta_k\omega_k/(\omega_k+\delta_k-x)^2<0$, $g_z(x)\to+\infty$ as $x\to-\infty$ and $g_z(x)\to-\infty$ as $x\uparrow d_{\min}$, there is exactly one root $x_0$ below $d_{\min}$, and every other eigenvalue is at least $d_{\min}$; so $x_0=\lambda_{\min}(\mathsf R)$. Since $\omega_k+\delta_k>0$ for every $k$, the point $0$ lies below $d_{\min}$, where $g_z$ decreases; hence $x_0\ge0$ if and only if $g_z(0)=\varepsilon'+\sum_k\beta_k\delta_k/(\omega_k+\delta_k)\ge0$. (c) By the Perron--Frobenius theorem $\rho(I-\mathsf R)$ is an eigenvalue of $I-\mathsf R$, and it is at least every real eigenvalue; as all eigenvalues are real, it is the largest, $1-x_0$. (d) For $A\ge0$ and a positive vector $y$, $\min_i(Ay)_i/y_i\le\rho(A)\le\max_i(Ay)_i/y_i$ (Collatz--Wielandt), so a positive eigenvector belongs to the spectral radius. Since the joint matrix is nonnegative, every agent's self-weight is at least $\sum_k\beta_k$ and $1-\omega_k-\delta_k\ge0$, so the reduced matrix $I-\mathsf R$ with $\upsilon_k=0$, whose first diagonal entry $1-\varepsilon-\sum_k\beta_k$ is at least an agent's self-weight minus $\sum_k\beta_k$, is nonnegative; it is irreducible, with a positive Perron vector $(y_0,y_1,\dots,y_{n_z})$. The vector $(y_0\one,y_1,\dots,y_{n_z})$ is an eigenvector of the joint matrix with eigenvalue $1-x_0$: an agent's row gives $(1-\varepsilon)y_0-\sum_k\beta_ky_0+\sum_k\beta_ky_k$, using $W_0\one=(1-\varepsilon)\one$, and a record's row gives $\sum_i\omega_{k,i}y_0+(1-\omega_k-\delta_k)y_k$, using $\sum_i\omega_{k,i}=\omega_k$; both equal the corresponding rows of $(I-\mathsf R)y$. It is positive, so $1-x_0$ is the spectral radius.
\end{proof}

\begin{proof}[Proof of \cref{prop:sca-records-memory}]
(a) Put $m_-=\min\{\varepsilon,\delta_k\}$ and $m_+=\max\{\varepsilon,\delta_k\}$, with $\varepsilon'=\varepsilon$. Every term of $g_z(m_-)=\varepsilon-m_-+\sum_k\beta_k(\delta_k-m_-)/(\omega_k+\delta_k-m_-)$ is nonnegative, so $x_0\ge m_-$. If $m_+<d_{\min}$, every term of $g_z(m_+)$ is nonpositive and $x_0\le m_+$; otherwise $x_0<d_{\min}\le m_+$. For one record with $\delta<\varepsilon$, $g_z(\delta)=\varepsilon-\delta>0$, and either $g_z(\varepsilon)=\beta(\delta-\varepsilon)/(\omega+\delta-\varepsilon)<0$ or $\varepsilon\ge\omega+\delta>x_0$; so $\delta<x_0<\varepsilon$, and the case $\delta>\varepsilon$ is symmetric. For $\varepsilon<d_{\min}$, $x_0<\varepsilon$ if and only if $g_z(\varepsilon)<0$, and $g_z(\varepsilon)$ is the sum in \eqref{eq:sca-memory-criterion}; for $\varepsilon\ge d_{\min}$, $x_0<d_{\min}\le\varepsilon$. The same signs show that the inequalities are strict when every $\delta_k<\varepsilon$ or every $\delta_k>\varepsilon$. When $I-\mathsf R\ge0$, $\tau=-1/\ln(1-x_0)$ by \cref{prop:sca-records}(c), and $x_0\ge m_-$ gives the cap on $\tau$.

(b) At $\delta=0$ the root satisfies $x_0=\varepsilon(\omega-x_0)/(\omega-x_0+\beta)$. Since $u\mapsto u/(u+\beta)$ is increasing and $\omega-x_0\le\omega$, $x_0\le\varepsilon\omega/(\omega+\beta)=\varepsilon/(1+\mathcal R)$; and $x_0\ge\varepsilon(\omega-x_0)/(\omega+\beta)$ rearranges to $x_0\ge\varepsilon\omega/(\omega+\beta+\varepsilon)$. Since $\partial g_z/\partial\delta=\beta\omega/(\omega+\delta-x)^2>0$ and $g_z$ decreases in $x$, the root increases with $\delta$.

(c) With $\varepsilon'=\varepsilon-\upsilon\beta$ and $\varepsilon<\omega+\delta$, $g_z(\varepsilon)=\beta[(\delta-\varepsilon)/(\omega+\delta-\varepsilon)-\upsilon]$, which is negative if and only if $(1-\upsilon)(\delta-\varepsilon)<\upsilon\omega$. If $\varepsilon\ge\omega+\delta$, then $\delta<\varepsilon$ and $x_0<\omega+\delta\le\varepsilon$, consistent with the criterion. Several records are treated in the same way.

(d) If $0\le A\le B$ entrywise, then $A^n\le B^n$, so $\Vert A^n\Vert_\infty\le\Vert B^n\Vert_\infty$ and Gelfand's formula gives $\rho(A)\le\rho(B)$; apply it to the matrix whose only nonzero entry is the record's self-weight. When reading displaces only leak, the agents keep their self-weight and gain the read weight, so the joint matrix dominates the agents-only matrix and $\rho\ge1-\varepsilon$, the case $\upsilon=1$ of (c).

(e) At $(\varepsilon,\delta)=0$, $g_z(0)=0$ and $\partial_xg_z(0)=-(1+\sum_k\mathcal R_k)\neq0$; $g_z$ is analytic near the origin, with $\partial_\varepsilon g_z=1$ and $\partial_{\delta_k}g_z=\beta_k\omega_k/(\omega_k+\delta_k-x)^2=\mathcal R_k$ there. The implicit function theorem gives the first formula of \eqref{eq:sca-records-tau}. Since $\varepsilon+\sum_k\mathcal R_k\delta_k\ge c_0(\varepsilon+\sum_k\delta_k)$ with $c_0>0$, the quadratic error is $o(\varepsilon_{\mathrm{rec}})$, and $\tau=-1/\ln(1-\varepsilon_{\mathrm{rec}})=\varepsilon_{\mathrm{rec}}^{-1}(1+O(\varepsilon_{\mathrm{rec}}))$.
\end{proof}

\begin{proof}[Proof of \cref{prop:sca-records-length}]
(a) For kernels symmetric under reflections $\hat{\mathsf W}(p)$ is real, and each of its entries is a Fourier sum of nonnegative real-space weights, so $|\hat{\mathsf W}(p)|\le\hat{\mathsf W}(0)$ entrywise and $\rho(\hat{\mathsf W}(p))\le\rho(\hat{\mathsf W}(0))=\rho(I-\mathsf R)<1$. The stationary covariance of mode $p$ is therefore the unique solution $\hat\Sigma(p)=\mathcal L_p^{-1}(Q_z)$ of $\hat\Sigma(p)=\hat{\mathsf W}(p)\hat\Sigma(p)\hat{\mathsf W}(p)^\top+Q_z$, where $\mathcal L_p(X)=X-\hat{\mathsf W}(p)X\hat{\mathsf W}(p)^\top$ is invertible. The agents' row has $\hat{\mathsf W}(p)_{00}=1-\varepsilon'-\sum_k\beta_k-a\sum_d2(1-\cos p_d)$ and $\hat{\mathsf W}(p)_{0k}=\beta_k\hat{\mathsf K}_k(p)$, with $\hat{\mathsf K}_k(p)=1-\ell_{r,k}^2|p|^2/2+o(|p|^2)$ by the finite second moments and the symmetry of the kernels; this gives the expansion with $\mathsf R_2$, and $\hat{\mathsf W}$ is twice continuously differentiable near $p=0$. The solution of the Lyapunov equation is a rational function of the entries of $\hat{\mathsf W}(p)$ without poles where $\rho(\hat{\mathsf W}(p))<1$, so $\hat\Sigma$ is twice continuously differentiable near the origin, with $\hat\Sigma(p)=\Sigma_0-|p|^2\Sigma_2+o(|p|^2)$; collecting the orders $|p|^0$ and $|p|^2$ gives \eqref{eq:sca-records-lyap}, each solution being unique because $\rho(\hat{\mathsf W}_0)<1$. The real-space transition matrix and the noise covariance are entrywise nonnegative, so the agents' covariance $\Sigma_{00}(x)$ is nonnegative and summable, with $\sum_x\Sigma_{00}(x)=(\Sigma_0)_{00}$. For a nonnegative summable sequence whose Fourier series is twice differentiable at the origin, Fatou's lemma applied to $(1-\cos(p\cdot x))/|p|^2$ shows that its second moments are finite, and dominated convergence then gives $\sum_x|x|^2\Sigma_{00}(x)=-\Delta_p\hat\Sigma_{00}(0)=2D(\Sigma_2)_{00}$. The definition $\xi_2^2=\sum_x|x|^2\Sigma_{00}(x)/(2D\sum_x\Sigma_{00}(x))$ gives the formula.

(b) Let $x_0=\varepsilon_{\mathrm{rec}}$ be the slow eigenvalue of $\mathsf R$, with right and left eigenvectors $r_0$ and $l_0$, $(r_0)_0=1$ and $l_0^\top r_0=1$. From the record rows of $\mathsf Rr_0=x_0r_0$ and the record columns of $l_0^\top\mathsf R=x_0l_0^\top$, $(r_0)_k=\omega_k/(\omega_k+\delta_k-x_0)$ and $(l_0)_k=(l_0)_0\beta_k/(\omega_k+\delta_k-x_0)$, and the normalization gives $(l_0)_0=(1+\sum_k\beta_k\omega_k/(\omega_k+\delta_k-x_0)^2)^{-1}$. Since only the first row of $\mathsf R_2$ is nonzero, $l_0^\top\mathsf R_2r_0=(l_0)_0(a+\sum_k\beta_k\ell_{r,k}^2(r_0)_k/2)$, which is \eqref{eq:sca-deff}; by first-order perturbation of a simple eigenvalue it is the $|p|^2$-coefficient of the slow eigenvalue of $I-\hat{\mathsf W}(p)=\mathsf R+|p|^2\mathsf R_2+o(|p|^2)$.

At zero leaks with $\upsilon_k=0$, $\mathsf R$ has zero row sums, so $0$ is an eigenvalue with right eigenvector $\one$, and $I-\mathsf R$ is nonnegative and irreducible (its graph is a star), with a positive diagonal entry $1-\sum_k\beta_k\ge2Da>0$. It is therefore primitive, and its other eigenvalues have modulus below one; equivalently, the other eigenvalues of $\mathsf R$ lie in a compact subset of $(0,2)$. By continuity this persists for small leaks, so $x_0$ is a simple isolated eigenvalue, and the spectral projector $\mathsf P_0=r_0l_0^\top$, the vectors $r_0$, $l_0$ and the fast part $M_f=(I-\mathsf R)(I-\mathsf P_0)$ depend continuously on the leaks. At zero leaks $\rho(M_f)<1$, so some norm makes $M_f$ a strict contraction there, and by continuity it remains one, with a common factor below one, near zero leaks; the powers of $M_f$ therefore decay geometrically, uniformly near zero leaks. At zero leaks $r_0=\one$ and $l_0=(1,\mathcal R_1,\dots,\mathcal R_{n_z})^\top/(1+\sum_k\mathcal R_k)$.

The map $\mathcal L(X)=X-\hat{\mathsf W}_0X\hat{\mathsf W}_0^\top$ commutes with the four block projections $X\mapsto\mathsf P_aX\mathsf P_b^\top$, $\mathsf P_a,\mathsf P_b\in\{\mathsf P_0,I-\mathsf P_0\}$, because $\hat{\mathsf W}_0=I-\mathsf R$ commutes with $\mathsf P_0$. On the slow-slow block it is multiplication by $1-(1-x_0)^2=x_0(2-x_0)$. On the other three blocks its inverse is a convergent series in $M_f$ (for example $Y\mapsto\sum_n(1-x_0)^nYM_f^{n\top}$ on the slow-fast block, where $0\le1-x_0\le1$), bounded uniformly near zero leaks by the geometric decay of the powers of $M_f$. With $\mathsf q_0=l_0^\top Q_zl_0$, the first equation of \eqref{eq:sca-records-lyap} therefore gives
\begin{equation*}
\Sigma_0=\frac{\mathsf q_0}{x_0(2-x_0)}r_0r_0^\top+B_0,\qquad\Sigma_0l_0=\frac{\mathsf q_0}{x_0(2-x_0)}r_0+B_0l_0,
\end{equation*}
with $B_0$ bounded. In the second equation the source $Y=\mathsf R_2\Sigma_0\hat{\mathsf W}_0^\top+\hat{\mathsf W}_0\Sigma_0\mathsf R_2^\top$ is $O(1/x_0)$, and its slow-slow block is $\mathsf P_0Y\mathsf P_0^\top=(1-x_0)r_0r_0^\top(l_0^\top\mathsf R_2\Sigma_0l_0+l_0^\top\Sigma_0\mathsf R_2^\top l_0)=r_0r_0^\top[2(1-x_0)\mathcal D_{\mathrm{eff}}\mathsf q_0/(x_0(2-x_0))+O(1)]$, using $l_0^\top\mathsf R_2r_0=\mathcal D_{\mathrm{eff}}$. Dividing the slow-slow block by $x_0(2-x_0)$ and bounding the other blocks,
\begin{equation*}
\Sigma_2=\frac{2(1-x_0)\mathcal D_{\mathrm{eff}}\mathsf q_0}{x_0^2(2-x_0)^2}r_0r_0^\top+O(1/x_0).
\end{equation*}
Since $(r_0)_0=1$ and $\mathsf q_0\to q/(1+\sum_k\mathcal R_k)^2+q_z\sum_k\mathcal R_k^2/(1+\sum_k\mathcal R_k)^2>0$, the agents' entries give
\begin{equation*}
\xi_2^2=\frac{(\Sigma_2)_{00}}{(\Sigma_0)_{00}}=\frac{2(1-x_0)}{2-x_0}\cdot\frac{\mathcal D_{\mathrm{eff}}}{x_0}\bigl(1+O(x_0)\bigr)=\frac{\mathcal D_{\mathrm{eff}}}{x_0}\bigl(1+O(x_0)\bigr),
\end{equation*}
the first form of \eqref{eq:sca-records-xi}, with constants depending on the fixed parameters. As the leaks vanish, $\omega_k/(\omega_k+\delta_k-x_0)\to1$ and $\beta_k\omega_k/(\omega_k+\delta_k-x_0)^2\to\mathcal R_k$, so $\mathcal D_{\mathrm{eff}}\to(a+\sum_k\beta_k\ell_{r,k}^2/2)/(1+\sum_k\mathcal R_k)>0$, and \cref{prop:sca-records-memory}(e) gives the second form.

(c) Let $\Pi$ be the joint precision of agents and records, diffusive by hypothesis, and $S=\Pi_{AA}-\Pi_{AZ}\Pi_{ZZ}^{-1}\Pi_{ZA}$ the agents' marginal precision. The proof of \cref{lem:col-kron}, with the agents as members and the records as outsiders, uses only the sign pattern of $\Pi$ and the Neumann series of $\Pi_{ZZ}^{-1}$, which converges in operator norm on $\Z^D$ because $\Pi\succeq\Vert\Sigma\Vert^{-1}I$ and the diagonal of $\Pi$ is bounded; so $S$ is diffusive, its leaks are at least those of $\Pi$ on the agents, and its couplings among agents are at least those of $\Pi$. $S$ is translation invariant, hence has a uniform leak, and $\hat S=1/\hat\Sigma_{00}$ is twice differentiable at the origin by (a); since its couplings are nonnegative, Fatou's lemma again gives them finite second moments, and \cref{prop:sca-lengths-relation}(a) shows that the length of (a) is the second-moment length of $S$. For a region $C$ of agents, $\Pi_{ZZ}^{-1}\succ0$ gives $S\preceq\Pi_{AA}$, so $\xout$ computed from $S$ is at most $\xout$ computed from $\Pi$; and since leaks and internal couplings only increase, the quadratic form that defines $\xin$ does not decrease on vectors supported on $C$. Hence the contrast computed from $S$ is at least the contrast computed from $\Pi$, which is at least $r_0$. \Cref{thm:col-contrast-size} with uniform leak, applied to $S$, gives $r_S(C)\le1+\xi_2^2/\Var_d(C)$, hence $\Var_d(C)\le\xi_2^2/(r_S(C)-1)\le\xi_2^2/(r_0-1)$.
\end{proof}

The long-range results rest on two lemmas about convolution powers of a probability $\mathsf J$ on $\Z^D\setminus\{0\}$; $\mathsf J^{\ast m}$ is the law of the sum $X_1+\cdots+X_m$ of independent steps with law $\mathsf J$.

\begin{lemma}[Union bound]\label{lem:sca-union}
If $\mathsf J(y)\le K_{\mathsf J}|y|^{-D-\sigma}$ for $y\neq0$, then $\mathsf J^{\ast m}(x)\le K_{\mathsf J}m^{1+D+\sigma}|x|^{-D-\sigma}$ for all $m\ge1$ and $x\neq0$.
\end{lemma}
% src: L4 Lemma 5.1
\begin{proof}
If $X_1+\cdots+X_m=x$ then $\max_l|X_l|\ge|x|/m$. By exchangeability, $\Pr(X_1+\cdots+X_m=x)\le m\Pr(X_1+\cdots+X_m=x,\ |X_m|\ge|x|/m)=m\sum_{|y|\ge|x|/m}\mathsf J(y)\mathsf J^{\ast(m-1)}(x-y)\le m\sup_{|y|\ge|x|/m}\mathsf J(y)$, which is at most $K_{\mathsf J}m^{1+D+\sigma}|x|^{-D-\sigma}$.
\end{proof}

\begin{lemma}[Single big jump]\label{lem:sca-big-jump}
If $\mathsf J(x)\sim\mathsf J_\infty|x|^{-D-\sigma}$ with $\sigma>0$, then $\mathsf J^{\ast m}(x)/\mathsf J(x)\to m$ as $|x|\to\infty$, for each fixed $m$.
\end{lemma}
% src: L4 Lemma 5.2
\begin{proof}
By induction it suffices to show that if $F$ and $G$ are probabilities with $F(x)\sim\ell_F\mathsf J(x)$ and $G(x)\sim\ell_G\mathsf J(x)$, both dominated by $K|x|^{-D-\sigma}$, then $(F\ast G)(x)\sim(\ell_F+\ell_G)\mathsf J(x)$. Fix $\delta_1\in(0,1/2)$ and split $\sum_yF(y)G(x-y)$ over $\mathcal A=\{|y|\le\delta_1|x|\}$, $\mathcal B=\{|x-y|\le\delta_1|x|\}$ and the rest. On $\mathcal A$, $G(x-y)/\mathsf J(x)$ lies in $\ell_G[(1+\delta_1)^{-D-\sigma},(1-\delta_1)^{-D-\sigma}](1+o(1))$ uniformly, and $F(\mathcal A)\to1$; $\mathcal B$ is symmetric with $\ell_F$. The rest is at most $\sup_{|z|>\delta_1|x|}G(z)\cdot F(\{|y|>\delta_1|x|\})=O(|x|^{-D-\sigma}|x|^{-\sigma})=o(\mathsf J(x))$. Let $\delta_1\to0$.
\end{proof}

\begin{proof}[Proof of \cref{prop:sca-longrange}]
$\Sigma=q\sum_nW^{2n}$ and $W^{2n}=\sum_{m=0}^{2n}\binom{2n}mc^{2n-m}a^m\mathsf J^{\ast m}$, so for $x\neq0$
\begin{equation*}
\frac{\Sigma(x)}{\mathsf J(x)}=q\sum_{n\ge1}\sum_{m=1}^{2n}\binom{2n}mc^{2n-m}a^m\frac{\mathsf J^{\ast m}(x)}{\mathsf J(x)}.
\end{equation*}
Each ratio tends to $m$ by \cref{lem:sca-big-jump}, and is dominated by $K'm^{1+D+\sigma}$ by \cref{lem:sca-union} together with $\mathsf J(x)\ge c_{\mathsf J}|x|^{-D-\sigma}$ for large $|x|$. Since $\sum_n\sum_m\binom{2n}mc^{2n-m}a^mm^{1+D+\sigma}\le\sum_n(2n)^{1+D+\sigma}\lambda^{2n}<\infty$, dominated convergence gives the limit $q\sum_n\sum_m\binom{2n}mc^{2n-m}a^mm=q\sum_n2na\lambda^{2n-1}=2qa\lambda/(1-\lambda^2)^2$. Without translation invariance, $\Vert\Sigma_{ij}\Vert\le q\sum_n\sum_k\Vert(W^n)_{ik}\Vert\Vert(W^n)_{jk}\Vert$, and $\Vert(W^n)_{ik}\Vert$ is dominated by the $n$-th convolution power of the symmetric majorant $c\one\{x=0\}+a\mathsf J(x)$, so the bound is the translation-invariant expression evaluated at $x_i-x_j$.
\end{proof}

\begin{proof}[Proof of \cref{prop:sca-fronts}]
The normalized influence is $\Pr(|Y_t|\ge r)$ for the walk $Y_t$ whose steps $X_s$ have law $(c/\lambda)\one\{x=0\}+(a/\lambda)\mathsf J(x)$. (a) $\Pr(|Y_t|\ge r)\le t\Pr(|X|\ge r/2)+\Pr(|\sum_sX_s\one\{|X_s|<r/2\}|\ge r)$. The tail of $\mathsf J$ gives $\Pr(|X|\ge r/2)=O(r^{-\sigma})$. The truncated steps are symmetric and independent, so Chebyshev's inequality bounds the second term by $t\E[|X|^2\one\{|X|<r/2\}]/r^2$, and the truncated second moment is $O(r^{2-\sigma})$ for $\sigma<2$, $O(\ln r)$ for $\sigma=2$ and $O(1)$ for $\sigma>2$. (b) Let $p_r=\Pr(|X|\ge2r)\ge c_3r^{-\sigma}$, and let $E_s$ be the event that $|X_s|\ge2r$, that $|X_{s'}|<2r$ for $s'\neq s$, and that $|\sum_{s'\neq s}X_{s'}|<r$. On $E_s$, $|Y_t|\ge r$, and the $E_s$ are disjoint. By independence and (a) applied to $t-1$ steps, once $r\ge c_2t^{1/\sigma}$ with $c_2$ large, $\Pr(E_s)\ge p_r(\tfrac12-tp_r)\ge p_r/4$, so $\Pr(|Y_t|\ge r)\ge tp_r/4\ge c_1tr^{-\sigma}$. (c) For $\sigma<2$, (a) gives $r_\theta(t)\le(C_1t/\theta)^{1/\sigma}$, and (b) at $r=Kt^{1/\sigma}$ with $c_2\le K<(c_1/\theta)^{1/\sigma}$ gives $r_\theta(t)\ge Kt^{1/\sigma}$, which is possible for $\theta<c_1c_2^{-\sigma}$. For $\sigma>2$ the steps have finite variance; (a) gives $r_\theta(t)=O(t^{1/2})$ and the central limit theorem gives the matching lower bound. For $\sigma=2$, (a) with the logarithm gives $r_\theta(t)=O((t\ln t)^{1/2})$. The front is at most linear exactly when $\sigma\ge1$. (d) $(\bar J^t)(x)\le\sum_m\binom tmc^{t-m}a^m\mathsf J^{\ast m}(x)$, and \cref{lem:sca-union} with $m\le t$ gives the bound; $(\bar J^t)_{ij}$ exceeds a level $\theta_1$ only if $t\ln g_{\mathrm{tot}}+(1+D+\sigma)\ln t\ge(D+\sigma)\ln|x_i-x_j|+\ln(\theta_1/K_{\mathsf J})$.
\end{proof}

\begin{proof}[Proof of \cref{thm:sca-percolation}]
All randomness other than $X_{\mathrm{src}}$, namely the initial states of the other agents and the fresh $\zeta_{i,t}$, is independent and can be absorbed into the channels, so the run is a Bayesian network with the single source $X_{0,0}$. Theorem~5 of \citet{polyanskiy2017strong} bounds the strong data-processing coefficient of the channel from $X_{0,0}$ to $X_{\mathcal S}$ by $\perc(X_{0,0}\to\mathcal S)$ for site percolation with each vertex open with the coefficient of its channel, and $I(X_{\mathrm{src}};X_{\mathcal S})\le s_{\mathrm{KL}}(P_{X_{\mathcal S}\mid X_{0,0}})I(X_{\mathrm{src}};X_{0,0})$ for the chain $X_{\mathrm{src}}\to X_{0,0}\to X_{\mathcal S}$ by the definition of the coefficient. Their Corollary~7 bounds the percolation probability by the sum over directed paths of the products of the open probabilities. A path from $X_{0,0}$ to $X_{i,t}$ makes $t$ steps, each a message step, open with probability at most $s_e$, or a memory step, open with probability at most $s_m$; it makes at least $d$ message steps; and there are at most $\binom tk\Delta^k$ paths with $k$ message steps. This gives the first bound. Summing over $t\ge k$ with $\sum_{t\ge k}\binom tky^{t-k}=(1-y)^{-k-1}$ and then over $k\ge d$ gives \eqref{eq:sca-perc-union}.
\end{proof}

\begin{proof}[Proof of \cref{thm:sca-trees}]
(a) Given $X_u$ the subtrees of the children are independent, and by induction each child's subtree is equivalent to an observation $X_c+N_c$ with $N_c\sim\Ncal(0,\mathcal V(c))$ independent of everything above $c$, since for Gaussian data the likelihood depends on the subtree only through one linear statistic. Then $(X_c+N_c)/\alpha_{e_c}=X_u+((1-\alpha_{e_c}^2)^{1/2}Z_c+N_c)/\alpha_{e_c}$ has noise variance $(1-\alpha_{e_c}^2+\mathcal V(c))/\alpha_{e_c}^2$, and independent observations add precisions. A childless vertex above level $n$ has no observed descendant, so its subtree is the uninformative observation $\mathcal V=\infty$ and contributes nothing to its parent's sum. The root's posterior variance is $\mathcal V/(1+\mathcal V)$, so $I=\frac12\ln(1+1/\mathcal V)$. (b) Substituting $\mathcal V(u)=\alpha^{2|u|}\hat{\mathcal V}(u)$, the recursion becomes $\hat{\mathcal V}(u)^{-1}=\sum_c((1-\alpha^2)\alpha^{-2|c|}+\hat{\mathcal V}(c))^{-1}$, the series and parallel laws, with terminal values $\hat{\mathcal V}=\sigma_{\mathrm{obs}}^2\alpha^{-2n}$ at level $n$. (c) By symmetry level $k$ contributes $b^k$ equal resistances $(1-\alpha^2)\alpha^{-2k}$ in parallel, and the levels are in series, so $\mathcal V_n=(1-\alpha^2)\sum_{k=1}^n(b\alpha^2)^{-k}$; expanding $\frac12\ln(1+1/\mathcal V_n)$ in the three regimes gives the limits. Because the data are Gaussian, the linear statistic is sufficient and linear estimation is optimal. (d) The resistance to the shorted level $n$ increases with $n$ to the resistance to infinity, and a locally finite network is transient exactly when that resistance is finite; by (b), $\lim_nI_n>0$ exactly then. The conductances $\alpha^{2|e|}/(1-\alpha^2)$ are proportional to $(\alpha^{-2})^{-|e|}$, which define the biased walk of \citet{lyons1990random} with bias $\alpha^{-2}$; by his theorem, as restated by \citet{lyons1992random}, that walk is transient if $\alpha^{-2}<\mathrm{br}(T)$ and recurrent if $\alpha^{-2}>\mathrm{br}(T)$.
\end{proof}

\begin{proof}[Proof of \cref{thm:sca-gaussian-percolation}]
Let $\mathfrak g(u)=1/(1+\mathcal V(u))$, the squared correlation of $X_u$ with its subtree statistic. We show $\mathfrak g(u)\le\mathfrak p(u)$ by induction from the observed level, where $\mathfrak g=1/(1+\sigma_{\mathrm{obs}}^2)\le1=\mathfrak p$, and from the childless vertices above it, where $\mathfrak g=0=\mathfrak p$. With $h_c=\alpha_{e_c}^2\mathfrak g(c)$, $1-\alpha_{e_c}^2+\mathcal V(c)=(1-h_c)/\mathfrak g(c)$, so the recursion of \cref{thm:sca-trees}(a) reads $\mathcal V(u)^{-1}=\mathcal T:=\sum_ch_c/(1-h_c)$ and $\mathfrak g(u)=\mathcal T/(1+\mathcal T)$, increasing in each $h_c$. By induction $h_c\le\bar h_c:=\alpha_{e_c}^2\mathfrak p(c)$, so $\mathfrak g(u)\le\bar{\mathcal T}/(1+\bar{\mathcal T})$ with $\bar{\mathcal T}=\sum_c\bar h_c/(1-\bar h_c)$. With $u_c=1/(1-\bar h_c)\ge1$, $1+\bar{\mathcal T}=1+\sum_c(u_c-1)\le\prod_cu_c$, hence $\bar{\mathcal T}/(1+\bar{\mathcal T})\le1-\prod_c(1-\bar h_c)=\mathfrak p(u)$. The observed level enters through the statistic $X_{\mathrm{root}}+N$ with $N\sim\Ncal(0,\mathcal V(\mathrm{root}))$ independent of $(X_{\mathrm{src}},X_{\mathrm{root}})$, which is sufficient, so the squared multiple correlation of $X_{\mathrm{src}}$ with the observations is $\mathfrak g(\mathrm{root})$ times that with $X_{\mathrm{root}}$. Finally $f(x)=-\frac12\ln(1-x)$ is convex with $f(0)=0$, so $f(\mathfrak gx)\le\mathfrak gf(x)$ for $\mathfrak g\in[0,1]$, and $\mathfrak g(\mathrm{root})\le\mathfrak p(\mathrm{root})$, the bond-percolation probability.
\end{proof}

\begin{proof}[Proof of \cref{prop:sca-aggregate}]
Let $Y=(e_j)_{d(0,j)\ge r}$. The squared top canonical correlation of $e_0$ and $Y$ is the largest eigenvalue of $\Sigma_{00}^{-1/2}\Sigma_{0Y}\Sigma_{YY}^{-1}\Sigma_{Y0}\Sigma_{00}^{-1/2}$, at most $\Vert\Sigma_{0Y}\Vert^2/(\lambda_{\min}(\Sigma_{YY})\lambda_{\min}(\Sigma_{00}))$. Interlacing gives $\lambda_{\min}(\Sigma_{YY})\ge\lambda_{\min}(\Sigma)\ge\sigma_{\mathrm e}$, and $\Vert\Sigma_{0Y}\Vert^2\le\sum_j\Vert\Sigma_{0j}\Vert^2\le B_\Sigma^2\sum_je^{-2m_\Sigma d(0,j)}$. If the spheres grow subexponentially, the tail sum decays exponentially in $r$.
\end{proof}

\begin{proof}[Proof of \cref{prop:sca-tree-diffusive}]
(a) The joint density is $p(x_o)\prod_cp(x_c\mid x_{\mathrm{parent}(c)})$ with conditional variance $1-\alpha^2$, so the precision has entries $-\alpha/(1-\alpha^2)$ on edges and diagonal entries $(1+b\alpha^2)/(1-\alpha^2)$ at internal vertices, $1+b\alpha^2/(1-\alpha^2)$ at the root and $1/(1-\alpha^2)$ at the leaves. The row sums follow, for example $(1+b\alpha^2-\alpha-b\alpha)/(1-\alpha^2)=(1-b\alpha)/(1+\alpha)$ at internal vertices. Off-diagonal entries are nonpositive, so the precision is diffusive exactly when every row sum is nonnegative. The leaf sums are positive, and the root sum is nonnegative exactly when $(b-1)\alpha\le1$. A tree of depth at least two has internal vertices, whose sums are nonnegative exactly when $b\alpha\le1$, a condition that implies the one at the root; a tree of depth one has none, and the root condition $(b-1)\alpha\le1$ alone decides. Since $\alpha<1$, $b\alpha^2>1$ forces $b\alpha>1$. (b) With $P=a^\Pi\mathcal A/(\varepsilon^\Pi+(b+1)a^\Pi)$, $\mathcal A$ the adjacency matrix, $P$ is the transition matrix of the simple random walk killed at rate $\varepsilon^\Pi/(\varepsilon^\Pi+(b+1)a^\Pi)$, and $\Sigma=(\varepsilon^\Pi+(b+1)a^\Pi)^{-1}\sum_mP^m$. The strong Markov property gives $\Sigma(o,v)=\Pr_o(\text{hit }v)\Sigma(v,v)$. Killing lowers hitting probabilities, and for the simple random walk on the $(b+1)$-regular tree the distance to $v$ decreases with probability $1/(b+1)$ and increases with probability $b/(b+1)$, so each unit decrease is ever achieved with probability $1/b$ and $\Pr_o(\text{hit }v)\le b^{-d(o,v)}$. By transitivity $\Sigma(v,v)=\Sigma(o,o)$, finite also at $\varepsilon^\Pi=0$ by transience. The squared canonical correlation with the sphere $S_n$ is at most $\Vert\Sigma_{oS_n}\Vert^2/(\lambda_{\min}(\Sigma_{S_nS_n})\Sigma(o,o))$, with $\Vert\Sigma_{oS_n}\Vert^2\le|S_n|b^{-2n}\Sigma(o,o)^2$, $|S_n|=(b+1)b^{n-1}$, and $\lambda_{\min}(\Sigma_{S_nS_n})\ge1/\lambda_{\max}(\Pi)\ge1/(\varepsilon^\Pi+2(b+1)a^\Pi)$ by Gershgorin's theorem.
\end{proof}

\begin{proof}[Proof of \cref{thm:sca-coherence-volume}]
(a) The variable $Z$ ranges over all nondegenerate square-integrable variables, on an extension of the probability space if necessary, so a variable independent of the readouts is admissible and $\theta^\ast\ge0$. Write $R=\mathrm{Corr}_C$ and $\underline v=\min_{p\in\mathbb P_C}p^\top Rp$. If $\underline v=0$, some $p\in\mathbb P_C$ has $p^\top X=0$ almost surely, so $\sum_ip_i\Corr(X_i,Z)=0$ for every $Z$, hence $\min_i\Corr(X_i,Z)\le0$ and $\theta^\ast=0$. Let $\underline v>0$. A variable $Z$ with $\min_i\Corr(X_i,Z)>0$ has a nonzero projection onto the linear span of the readouts, and the projection keeps every $\Cov(X_i,Z)$ and does not increase $\Var Z$, so it does not decrease any positive correlation; variables with nonpositive minimum do not affect a positive supremum. We may therefore take $Z=c^\top X$ with $c$ in the range of $R$. Normalizing $c^\top Rc=1$, $\Corr(X_i,Z)=(Rc)_i$, and the supremum of $\min_i(Rc)_i$ over the sphere $c^\top Rc=1$ agrees with the maximum over the ball $c^\top Rc\le1$ whenever the latter is positive, since scaling $c$ up increases a positive minimum. The maximum over the ball is $\max_{c^\top Rc\le1}\min_{p\in\mathbb P_C}p^\top Rc$. Both sets are compact and convex and the objective is bilinear, so von Neumann's minimax theorem gives $\min_p\max_cp^\top Rc=\min_p(p^\top Rp)^{1/2}$, by the Cauchy--Schwarz inequality in the inner product of $R$. At the saddle point $c^\ast=p^\ast/(p^{\ast\top}Rp^\ast)^{1/2}$, and the optimality conditions of the quadratic program on the simplex, $(Rp^\ast)_i\ge p^{\ast\top}Rp^\ast$ with equality on the support of $p^\ast$, show that $\min_i(Rc^\ast)_i=(p^{\ast\top}Rp^\ast)^{1/2}=\underline v^{1/2}>0$, so $\theta^\ast=\underline v^{1/2}$ is positive and attained by $c^{\ast\top}X$. (b) By (a), $\theta^{\ast2}=\underline v$ in both cases; take $p$ uniform, then use the Rayleigh quotient and Gershgorin's theorem. (c) Apply (b) to the set of members with $\Corr(X_i,Z)\ge\theta$, whose best correlational coherence is at least $\theta$ and whose susceptibility is at most $\chi_C$. (d) With $|\mathbf w|=1$, $Z=\mathbf w^\top X$ has $\Corr(X_i,Z)=(R\mathbf w)_i/(\mathbf w^\top R\mathbf w)^{1/2}=\lambda_{\max}^{1/2}w_i$.
\end{proof}

\begin{proof}[Proof of \cref{cor:sca-size}]
(a) From the series $\Sigma=\sum_n\Phi Q^\zeta\Phi^\top$ with block-diagonal noise, $\sum_j\Vert\Sigma_{ij}\Vert\le q\sum_n\sum_k\Vert\Phi_{ik}\Vert\sum_j\Vert\Phi_{jk}\Vert\le q\sum_n\lambda_r^n\lambda_c^n$, by the row and column inequalities. Readout correlations satisfy $|(\mathrm{Corr}_C)_{ij}|\le\Vert\Sigma_{ij}\Vert/(\sigma_i\sigma_j)$ with $\sigma_i^2=\mathbf v_i^\top\Sigma_{ii}\mathbf v_i\ge q_{\min}$, so $\chi_C\le q/(q_{\min}(1-\lambda_r\lambda_c))$, and \cref{thm:sca-coherence-volume}(b) concludes. (b) $\lambda_{\max}(\Sigma)\le\sum_n\Vert W^nQ^\zeta W^{n\top}\Vert\le K_A^2q/(1-f^2)$. The readout map has orthonormal columns in disjoint blocks and the standardization divides by $\sigma_i\ge q_{\min}^{1/2}$, so $\lambda_{\max}(\mathrm{Corr}_C)\le\lambda_{\max}(\Sigma)/q_{\min}$; and $1/(1-f^2)=\tau_f(1+O(1-f))/2$. Since $\Vert W^k\Vert\le K_Af^k$ forces $\rho(W)\le f$, $\tau_f\ge\tau$. (c) $W^n=(\lambda^n/N)\one\one^\top$ for $n\ge1$, so $\Sigma=qI+qb_N\one\one^\top$ with $b_N=\lambda^2/((1-\lambda^2)N)$, all correlations equal $b_N/(1+b_N)$, and by symmetry and convexity the uniform vector minimizes $p^\top\mathrm{Corr}_Vp$. Then $N\theta^{\ast2}=(1+Nb_N)/(1+b_N)$, which is also $\lambda_{\max}(\mathrm{Corr}_V)$, and $1+Nb_N=1/(1-\lambda^2)$. (d) $W\ge0$ gives $\Sigma\ge0$, so $\chi=\sum_x\Sigma(x)/\Sigma(0)=\hat\Sigma(0)/\Sigma(0)=q/((1-\lambda^2)\Sigma(0))$; insert \cref{thm:sca-oz}(f) and (g), with $1-\lambda^2\sim2\varepsilon$ and $\xi_{\mathrm{OZ}}\sim(a/\varepsilon)^{1/2}$. For $D\ge3$, $q/(1-\lambda^2)=q\xi_2^2/(2a\lambda)$. (e) The hub's error is an autoregression independent of the leaves, $e_h(t+1)=c_he_h(t)+\zeta_h(t)$, and each leaf has $e_\ell(t)=\sum_{n\ge0}c^n[ae_h(t-1-n)+\zeta_\ell(t-1-n)]$. For two leaves, $\Cov(e_\ell,e_{\ell'})=a^2\sum_{n,n'}c^{n+n'}\Cov(e_h(t-1-n),e_h(t-1-n'))>0$, the same for every pair and independent of $N$, as is $\Var(e_\ell)$; so all leaf-leaf correlations equal some $\theta_\star>0$. As in (c), a set of $N$ leaves has $\theta^{\ast2}=\theta_\star+(1-\theta_\star)/N>\theta^2$ when $\theta^2<\theta_\star$, for every $N$. The hub's column of $W$ has Euclidean norm $(c_h^2+Na^2)^{1/2}$, so $\Vert W\Vert\ge(c_h^2+Na^2)^{1/2}$, and the case $k=1$ of \textup{(C$_\rho$)} gives $K_Af\ge\Vert W\Vert$; since $f<1$, $K_A>a\sqrt N$, so the family satisfies neither hypothesis uniformly in $N$.
\end{proof}

\begin{proof}[Proof of \cref{thm:sca-alternating}]
(a) Each $P_{V_j}$ fixes $V_\cap$ pointwise and is self-adjoint, so $P_{V_j}P_\cap=P_\cap=P_\cap P_{V_j}$, hence $A_{\mathrm{cyc}}P_\cap=P_\cap A_{\mathrm{cyc}}=P_\cap$. Expanding with these relations, $(A_{\mathrm{cyc}}-P_\cap)^n=A_{\mathrm{cyc}}^n-P_\cap$, and $A_{\mathrm{cyc}}-P_\cap=A_{\mathrm{cyc}}(I-P_\cap)$. If $x\perp V_\cap$, $|x|=1$ and $|A_{\mathrm{cyc}}x|=1$, then each projection in turn preserves the norm and therefore fixes the vector, so $x\in\bigcap_jV_j=V_\cap$, a contradiction; since $A_{\mathrm{cyc}}$ maps $V_\cap^\perp$ into itself and the unit sphere of $V_\cap^\perp$ is compact, $c_{\mathrm{cyc}}<1$. (b) $V_\cap^\perp$ is invariant under every $P_{V_j}$, every product fixes $P_\cap x_0$, and $x_n-P_\cap x_0$ is the product applied to $x_0-P_\cap x_0\in V_\cap^\perp$. Every complete block of $p_0$ hops is a word containing every index, whose norm on $V_\cap^\perp$ is less than one by the argument of (a), and there are finitely many such words; the remaining hops have norm at most one. (c) Each factor has norm at most $\bar a$.
\end{proof}

\begin{proof}[Proof of \cref{thm:sca-h4}]
(a) The independence statement is \cref{thm:sca-lightcone}(b); the radius bound is \cref{thm:sca-no-moment}(b); \cref{def:sca-lengths} defines its lengths from the stationary covariance under \textup{(M)} or from the covariance of the process started in the infinite past under the uniform form, so it defines no length when both forms fail; the failure of the stationary form is by definition the absence of a stationary law with finite second moments, and \cref{def:col-collective} requires \textup{(M)}. (b) Under \textup{(B)}, standardized blocks satisfy $\Vert R_{ij}\Vert\le\Vert\Sigma_{ij}\Vert/q_{\min}$. \Cref{thm:sca-clustering} gives the first bound, \cref{prop:sca-spectral} the second, and \cref{thm:sca-diffusive}, whose prefactor is at most $2q\varpi_{\max}/\varpi_{\min}$, gives decay rate at least $\vartheta\sqrt{2\ln(1/\lambda)}$ for every $\vartheta<1$, hence the third. The lattice statements are \cref{thm:sca-oz}(c),(g) and \cref{prop:sca-lengths-relation}(c). (c) is \cref{prop:sca-lengths-relation}(b), \cref{prop:sca-longrange} and \cref{prop:sca-fronts}(c). (d) is \cref{prop:sca-persistence}, \cref{prop:sca-diffusive-scaling}(a), \cref{prop:sca-records}(c), \cref{prop:sca-records-memory}(a),(e) and \cref{prop:sca-records-length}(b). (e) is \cref{thm:col-contrast-size} with uniform leak, \cref{cor:col-lstar}, \cref{prop:sca-records-length}(c), \cref{thm:sca-coherence-volume}(c) and \cref{cor:sca-size}(a),(b),(e).
\end{proof}
